\documentclass[11pt,letterpaper]{article}
\usepackage[T1]{fontenc}
\usepackage{lmodern}
\usepackage[margin=0.88in]{geometry}
\usepackage{amsmath,amssymb,amsthm,mathtools}
\usepackage{booktabs,longtable,array,float}
\usepackage{microtype,xcolor,enumitem}
\usepackage{fancyhdr}
\usepackage{hyperref,xurl}
\hypersetup{colorlinks=true,linkcolor=blue!45!black,urlcolor=blue!45!black,citecolor=blue!45!black,
 pdftitle={Reciprocal pooling and integral projections for Diophantine consistency},
 pdfauthor={AI-assisted research prepared for Vladimir Reshetnikov}}
\pagestyle{fancy}\fancyhf{}
\fancyhead[L]{\small Diophantine consistency: reciprocal pooling}
\fancyhead[R]{\small 25 October 2026}
\fancyfoot[C]{\thepage}
\setlength{\headheight}{14pt}
\setlength{\parskip}{4pt}\setlength{\parindent}{0pt}
\setlength{\emergencystretch}{2em}
\newtheorem{theorem}{Theorem}[section]
\newtheorem{lemma}[theorem]{Lemma}
\newtheorem{corollary}[theorem]{Corollary}
\theoremstyle{definition}\newtheorem{definition}[theorem]{Definition}
\theoremstyle{remark}\newtheorem{remark}[theorem]{Remark}
\newcommand{\Z}{\mathbb Z}\newcommand{\N}{\mathbb N}
\newcommand{\Con}{\operatorname{Con}}
\newcommand{\ZFC}{\mathrm{ZFC}}
\newcommand{\Refut}{\operatorname{Ref}}
\newcommand{\sha}[1]{{\footnotesize\texttt{\seqsplit{#1}}}}
\newcommand{\code}[1]{\texttt{\detokenize{#1}}}
\newcommand{\Pq}{P_{218}}\newcommand{\Pd}{P_{215}}
% Frozen dimensions; full metric records are in equations/.
\newcommand{\PrincipalVariables}{218}
\newcommand{\CompanionVariables}{215}

\begin{document}
\begin{titlepage}
\vspace*{0.65in}
{\color{blue!45!black}\Large ProveIt research submission\par}
\vspace{0.35in}
{\Huge\bfseries Reciprocal pooling\\[8pt]and integral projections\par}
\vspace{0.25in}
{\Large Two explicit Diophantine equations for the syntactic consistency of ZFC\par}
\vspace{0.45in}
{\large A quartic in 218 integer unknowns\par}
{\large A degree-ten companion in 215 integer unknowns\par}
\vspace{0.6in}
Prepared for Vladimir Reshetnikov\par
AI-assisted mathematical research, 25 October 2026\par
\vfill
\begin{minipage}{0.96\textwidth}
\textbf{Result and scope.} This submission gives two complete, fixed
integer-coefficient polynomials, with literal residual lists, expanded
coefficients, ordinary mathematical proofs, and independent exact checkers.
Composing the new reductions with the retained direct C222/refutation
calibration proves in ZFC that each polynomial has no integer zero if and
only if $\Con(\ZFC)$.

The new arithmetic idea merges power certificates at \emph{different}
bases through integrality of a rational sum of squared errors. A common
sign and a strict size bound replace separate divisibility tests. The
quartic combines the newest same-base parent with seven exact integer
coordinate eliminations.

No global minimum, literature-wide priority, proof of consistency, or new
end-to-end proof-assistant certificate is claimed. The polynomial data
are complete; the logical calibration is an identified inherited theorem.
\end{minipage}
\end{titlepage}

\begin{abstract}
Starting from the repository's direct 225-integer quartic, we construct a
quartic in 218 unrestricted integer variables and 155 quadratic
residuals. Its expanded polynomial has 34,984 nonzero monomials and
coefficient height 113,304. Five common-Hessian directions and two affine
relations provide exact integral projections; a protected change of
residual basis lowers coefficient height. The larger support is an
explicit cost of this reduction.

We then prove a distinct-base reciprocal-pooling theorem. Fourteen power
requests at the consecutive bases $q$ and $q+1$ can share one integer
quotient, replacing four quadratic residuals and four private coordinates
by one quintic residual and one coordinate. The resulting single equation
has 215 unrestricted integer variables, exact degree ten, 152 residuals,
99,914 collected monomials, and the same coefficient height 113,304.
Its equivalence with the parent is proved in both directions, including
negative indices, signed Pell parameters, unit bases, and the bootstrap
that makes both bases nonzero. No coprimality between the two moduli is
required. Both polynomials encode the direct numerical proof calculus,
without K215, Jones's universal polynomial, or Sun's reduction.
\end{abstract}

\setcounter{tocdepth}{1}
\tableofcontents
\newpage

\section{The equations and the exact claim}

\subsection{A completely specified quartic}

All 218 coordinates are independent unrestricted integers. Define
\begin{equation}\label{eq:principal}
 \boxed{\Pq(x_0,\ldots,x_{217})=\sum_{i=0}^{154}Q_i(x_0,\ldots,x_{217})^2=0.}
\end{equation}
Appendix~\ref{app:residuals} prints every coefficient of every $Q_i$.
Each $Q_i$ has degree at most two. The file
\code{equations/quartic218/equation.json} gives the same list in an exact
sparse format; its companion \code{expanded_polynomial.json} gives all
34,984 collected terms of the single quartic. Thus the summation in
\eqref{eq:principal} is a finite polynomial abbreviation, not a
variable-length computation or an unspecified universal-polynomial index.

For the degree-ten companion, let $v=(v_0,\ldots,v_{214})\in\Z^{215}$ and put
\[
 \widehat Q_i(v)=Q_i(v_0,\ldots,v_{213},0,0,0,0),\qquad
 q=v_{95}-1,\quad A=v_{135}+1.
\]
Define the following polynomial abbreviations:
\begin{align}
 T_0&=2Aq-q^2-1,&T_1&=2A(q+1)-(q+1)^2-1,\label{eq:Tactual}\\
 E_0&=\widehat Q_{154}+q\widehat Q_{153},&
 E_1&=\widehat Q_{152}+(q+1)\widehat Q_{151},\label{eq:Eactual}\\
 \mathcal R&=E_0T_1+E_1T_0-v_{214}T_0T_1.\label{eq:Ractual}
\end{align}
The second equation is
\begin{equation}\label{eq:companion}
 \boxed{\Pd(v_0,\ldots,v_{214})=
   \sum_{i=0}^{150}\widehat Q_i(v)^2+\mathcal R(v)^2=0.}
\end{equation}
Equations~\eqref{eq:Tactual}--\eqref{eq:companion}, together with the
literal appendix, specify the second polynomial completely. Its
coefficient files are under \code{equations/decic215/}. There is no
division, exponent variable, factorial, binomial primitive, or hidden
coefficient parameter in either final polynomial.

\begin{theorem}[Consistency equivalence]\label{thm:main}
For the effective ZFC axiom predicate and finite proof system used in
the retained C222 calibration,
\begin{align*}
 \ZFC&\vdash
 \bigl(\exists x\in\Z^{218}\;\Pq(x)=0\bigr)
       \ \longleftrightarrow\ \exists p\,\Refut_{\ZFC}(p),\\
 \ZFC&\vdash
 \bigl(\exists v\in\Z^{215}\;\Pd(v)=0\bigr)
       \ \longleftrightarrow\ \exists p\,\Refut_{\ZFC}(p).
\end{align*}
Consequently, ZFC proves the absence of integer solutions of either
equation equivalent to $\Con(\ZFC)$.
\end{theorem}

\textbf{Meaning of the proof claim.} The new reductions below are
ordinary mathematical proofs formalizable in ZFC. Their composition uses
the inherited ordinary arithmetic proof of the 225-variable parent and
the named finite syntactic translations. The archive preserves those
sources and interfaces. It does not contain a newly checked, complete
Lean/Rocq derivation of Theorem~\ref{thm:main}, or a literal encoded ZFC
proof object. This is a distinction about the form of verification, not
an additional assumption that ZFC is consistent.

\subsection{Measured comparison}

\begin{table}[H]
\centering\small
\begin{tabular}{lrrrrr}
\toprule
Equation & Variables & Residuals & Degree & Terms & Height\\
\midrule
Pinned parent $P_{225}$ &225&162&4&7,867&128,446\\
Radical projection $P_{220}$ &220&157&4&7,790&142,238\\
\textbf{Principal $\Pq$} &\textbf{218}&\textbf{155}&\textbf{4}&\textbf{34,984}&\textbf{113,304}\\
\textbf{Companion $\Pd$} &\textbf{215}&\textbf{152}&\textbf{10}&\textbf{99,914}&\textbf{113,304}\\
\bottomrule
\end{tabular}
\caption{All domains are unrestricted integers. Terms and height refer
to the fully expanded, collected single polynomial.}
\label{tab:metrics}
\end{table}

Relative to its actual 225-variable parent, the principal saves seven
variables and seven residuals, retains degree four, and lowers height by
15,142, about 11.8\%. Its support increases by 27,117 monomials. The
215-variable companion saves another three variables and three residuals
at the cost of degree ten and a further 64,930 monomials. The maximum
absolute coefficient in an individual residual is 240 in both new
equations, versus 256 in the parent.

These are improvements in the direct, modest-coefficient C222 branch.
They are not absolute fewest-variable records. The inspected repository
already has a 37-integer quartic on a different coefficient-heavy branch,
and low-variable high-degree equations; its maintained display uses 15
positive integers, degree 157, and 66 integer-leaf binary operations.
Those measures describe different compromises~\cite{catalog,current}.
We do not replace a Pareto comparison by a single unsupported claim of
``the simplest equation.''

\section{The source and the logical connection}

\subsection{Pinned source}

The comparison snapshot is ProveIt commit
\begin{center}\small\texttt{953fb4d82d4229f4efc5f9db4cc2bf38ae329bb4}.\end{center}
The immediate mathematical parent is package
\code{158-pooled-powers}, whose principal residual file has SHA-256
\begin{center}\footnotesize\texttt{fee84c408828aba3e7ceebe9a37f58509327063aa7fd078fd49c80adf3af4e56}.\end{center}
That package supplies the 225-variable direct quartic and a complete
ordinary proof of its arithmetic interpretation~\cite{parent158}. The
generic quadratic-projection ideas and discovery utilities in package
159 supply a second input~\cite{projection159}. Applying those ideas to
the new 225-variable parent is a new instance; merely repeating a
projection lemma is not presented as a novel theorem.

The source route is
\[
\begin{aligned}
 \text{integer polynomial zero}
 &\ \longleftrightarrow\ \text{finite C222 refutation}\\
 &\ \longleftrightarrow\ \text{finite ZF refutation}\\
 &\ \longleftrightarrow\ \text{finite ZFC refutation}.
\end{aligned}
\]
The last two bridges are the retained interfaces
\begin{quote}\small
\code{KEndpoint.c222_iff_zf_refutation}\\
\code{ConstructibleChoice.zf_refutation_iff_zfc_refutation}.
\end{quote}
The source proof also retains its square-shell syntax recoding and
signed-Pell representation. All these are dependencies of the result;
the present work does not claim to have independently rebuilt the full
transitive formal development.

\subsection{What the arithmetic source certifies}

The C222 route encodes a finite derivation directly. A sufficiently large
finite horizon, factorial weights, and CRT residues encode the local
proof-rule tuples and the evolution of pending formulas. Exact powers
provide binomial digit and range certificates. The signed Pell framework
enables all final variables to range over $\Z$ while recovering the
semantic signs needed by the proof interpretation.

For the change made here, the crucial retained consequences are precise:
\begin{enumerate}[label=(\roman*)]
\item There is an integer $W$ with
 $W=2+\sum_{j=1}^{36}a_j^2$ for the listed affine frame atoms. The
 horizon atom $t=16+t_0^2$ implies $W\ge258$.
\item Before the two pooled arithmetic blocks are decoded, the common
 Pell parameter satisfies $|A|>W^W$ and the retained Pell pairs have
 their canonical signed indices.
\item Every actual base, requested output, and requested index of the
 fourteen pooled requests has square strictly below $W$.
\item Three singleton requests are retained: $(F^2+1)^t=V$, $B^x=N$,
 and $(BH)^x=q$. The normalized ceiling/remainder argument yields
 $|B|\ge2$; the multiplier $H$ is odd and nonzero.
\end{enumerate}
The parent proves these facts in an orientation-first, noncircular
order. Its proof of the common frame is not replaced by a numerical
experiment. The exact changes of coordinates below transport the entire
residual system, including these rows and their witnesses.

\subsection{Why this is an ordinary Diophantine equation}

The powers $W^W$ and $n^k$ appear in the \emph{proof} of the arithmetic
interpretation. They are not operations in the final polynomial. The
finite quadratic residuals encode the corresponding witnesses, and
\eqref{eq:companion} adds only a quintic polynomial. This separation is
essential: a short exponential-Diophantine formula would not by itself
answer the question posed here. The classical reduction literature also
distinguishes the number of unknowns, degree, and witness domain rather
than treating them as interchangeable measures~\cite{mr75,bayer2025}.

\section{Seven integral eliminations at degree four}

\subsection{Common-Hessian directions}

\begin{lemma}[Quadratic projection along a common radical]\label{lem:radical}
Let $R_0,\ldots,R_{m-1}\in\Z[x_0,\ldots,x_{n-1}]$ have degree at most
two. Suppose $d\in\Z^n$ lies in the kernel of every Hessian, $d_p=1$,
and for some row $r$ the integer
$a_r=\sum_jd_j[x_j]R_r$ is $1$ or $-1$. Put $a_i=\sum_jd_j[x_j]R_i$.
For $y_p=0$, define
\[
 Q_i(y)=R_i(y)-a_i a_rR_r(y)\quad(i\ne r).
\]
Then the old and new integer zero sets are in explicit bijection. One
variable and one residual disappear, and all remaining residuals are
quadratic.
\end{lemma}
\begin{proof}
The common-Hessian condition gives the exact polynomial identity
$R_i(y+sd)=R_i(y)+sa_i$. Set $s=-a_rR_r(y)$; this makes row $r$ zero
and gives $R_i(y+sd)=Q_i(y)$ for every other row. Conversely, for an old
zero $x$, put $y=x-dx_p$. Then $y_p=0$ and the pivot equation gives
$x_p=-a_rR_r(y)$. These two integral maps are inverse on the zero sets.
No division by a nonunit occurs. The formula for $Q_i$ is a constant row
combination, so its degree stays at most two.
\end{proof}

Five applications give the 220-variable intermediate. The following
indices refer to the input at each step; they are not all indices in the
original 225-variable system.
\begin{center}\small
\begin{tabular}{rrrrl}
\toprule
Dimension & Pivot $p$ & Deleted row $r$ & $a_r$ & Direction $d$\\
\midrule
225&44&29&1&$e_{44}-e_{45}$\\
224&60&46&$-1$&$e_{60}-e_{61}-2e_{62}$\\
223&64&49&1&$e_{64}-e_{62}$\\
222&31&16&1&$e_{31}-e_{32}-2e_{33}$\\
221&25&10&1&$e_{25}-e_{27}-2e_{28}$\\
\bottomrule
\end{tabular}
\end{center}
The frozen certificates include the complete quadratic restoration
polynomial at each step. The independent checker substitutes those maps
into every residual and checks both inverse identities. It does not
infer equivalence from matching dimensions or sampled evaluations.

\subsection{An integral chart of an affine hyperplane}

In the 220-variable system, the difference of rows 19 and 6 is
\begin{equation}\label{eq:hyperplane}
 4+12x_9+16x_{10}+4x_{20}-x_{21}+x_{34}=0.
\end{equation}
A direct substitution can make coefficients large. Instead use the
following integral parametrization, with four new coordinates $u,v,w,z$:
\begin{align*}
 x_9&=v+2w,&x_{10}&=-v-w,&x_{20}&=-1+v-2w+z,\\
 x_{21}&=u+2z,&x_{34}&=u-2z.
\end{align*}
It satisfies \eqref{eq:hyperplane} identically. Every integer point on
the hyperplane has the integral inverse
\begin{align*}
 v&=-x_9-2x_{10},&w&=x_9+x_{10},\\
 z&=1+3x_9+4x_{10}+x_{20},&u&=x_{21}-2z.
\end{align*}
The equation of the hyperplane ensures the restored $x_{34}$ is correct.
Substitute the chart in every row and omit row 6. Row 19 remains, so their
known difference restores the omitted row. This gives 219 variables and
156 residuals. The maps are explicit; no numerical lattice-reduction
heuristic is a premise of the proof.

\subsection{The final affine elimination}

In the resulting 219-variable chart, the surviving affine row is
\begin{align*}
0={}&4-4y_{10}-y_{20}-4y_{21}
 -2(y_{23}+y_{27}+y_{29}+y_{31}+y_{33})\\
&-4y_{35}-2(y_{36}+y_{39}+y_{42})-4y_{61}-2y_{62}+y_{94}.
\end{align*}
Its coefficient at $y_{94}$ is one. Solve for $y_{94}$, substitute in
every other row, and delete that row and coordinate. The inverse inserts
the displayed affine expression. This gives 218 variables and 155
quadratic residuals over exactly the same integer domain.

\subsection{A protected change of residual basis}

The raw affine projection has a large coefficient height. We apply 669
literal operations of the form $R_i\leftarrow R_i+cR_j$, $c\in\Z$,
$i\ne j$. Each is invertible by subtraction, so the common zero set is
unchanged. The certificate lists every operation in order and the
independent checker also replays the inverses in reverse order.

The last four rows, 151--154, are protected: they are neither changed nor
used as source rows in this selected basis transformation. This preserves
the two old power pools literally, which simplifies the next reduction.
The resulting sum of squares is the principal \eqref{eq:principal}.
The polynomial itself changes when the residual basis changes; the
integer zero set is preserved. The improved height therefore must be
measured from the newly expanded polynomial, as it is in
Table~\ref{tab:metrics}.

\section{Reciprocal pooling across distinct bases}

\subsection{The canonical Pell congruence}

For an integer $A$ with $|A|\ge2$, define the signed canonical Pell
coordinates by
\[
 (X_k,Y_k)=(A,1)^k\quad(k\in\Z)
\]
in the group with multiplication
\[
 (X,Y)(X',Y')=(XX'+(A^2-1)YY',\;XY'+YX').
\]
The identity is $(1,0)$ and $(A,1)^{-1}=(A,-1)$.
Thus every coordinate is an integer and
$X_k^2-(A^2-1)Y_k^2=1$.

Let $n\ne0$ and set $T=2An-n^2-1$. Modulo $T$,
$(n-A)^2\equiv A^2-1$. Therefore the map
$(X,Y)\mapsto X+(n-A)Y$ is multiplicative and sends $(A,1)$ to $n$.
Also $T\equiv-1\pmod n$, so $n$ is invertible modulo $T$. It follows for
all signed indices that
\begin{equation}\label{eq:pellcongruence}
 X_k+(n-A)Y_k\equiv n^k\pmod T.
\end{equation}
Negative powers in this congruence mean powers of the modular inverse.
They will be cleared before any size estimate is used.

\subsection{Statement of the new compression lemma}

\begin{theorem}[Two-base reciprocal pooling]\label{thm:pool}
Let $W\in\Z$, $W\ge258$, $A\in\Z$, and $|A|>W^W$. Let
$n_0=q$, $n_1=q+1$ be nonzero integers with $n_0^2<W$ and $n_1^2<W$. Partition at most 15 requests
into two sets $I_0,I_1$. For each request $i\in I_j$, let
$m_i,k_i\in\Z$ satisfy
\[
 n_j^2<W,\qquad m_i^2<W,\qquad k_i^2<W,
\]
and let $(X_i,Y_i)=(X_{k_i}(A),Y_{k_i}(A))$ be its canonical signed
Pell coordinates. Set $Z_i=AY_i$, and define
\begin{align*}
 T_j&=2An_j-n_j^2-1,\\
 H_j&=\sum_{i\in I_j}Y_i(m_i-X_i),\\
 S_j&=\sum_{i\in I_j}\bigl[(m_i-X_i+Z_i)^2-Y_i^2\bigr],\\
 E_j&=S_j-2n_jH_j.
\end{align*}
Then the following are equivalent:
\begin{enumerate}[label=(\alph*)]
\item Every request is the exact power $m_i=n_j^{k_i}$; when $|n_j|\ge2$,
 every $k_i$ is nonnegative. For a unit base the exact signed unit power
 is meant, without an index-sign conclusion.
\item There exists $z\in\Z$ such that
\begin{equation}\label{eq:RP}
 \boxed{E_0T_1+E_1T_0-zT_0T_1=0.}
\end{equation}
\end{enumerate}
The moduli $T_0,T_1$ need not be coprime.
\end{theorem}

\subsection{Proof: integrality, a common sign, and a strict bound}

Put $b=\lfloor\sqrt{W-1}\rfloor$. Then $b\ge16$ and every bounded
base, output, and index has absolute value at most $b$. First,
\begin{equation}\label{eq:Tbound}
 |T_j|\ge2|A||n_j|-n_j^2-1\ge2|A|-W>W^W.
\end{equation}
The sign of $T_j$ is the sign of $An_j$. Nonzero consecutive integers
have the same sign, so $T_0$ and $T_1$ have the same sign.

Write $\delta_i=m_i-X_i+Z_i-n_jY_i$. Direct expansion gives
\begin{equation}\label{eq:squareidentity}
 \sum_{i\in I_j}\delta_i^2=E_j-T_j\sum_{i\in I_j}Y_i^2.
\end{equation}
Consequently \eqref{eq:RP} says that
\begin{equation}\label{eq:rawfraction}
 \sum_{j=0}^1\frac{\sum_{i\in I_j}\delta_i^2}{T_j}\in\Z.
\end{equation}
The quantities $\delta_i$ may be very large. Bounding them directly would
be incorrect. Instead define the nonzero integer
\[
 C=(n_0n_1)^b
\]
and the \emph{small cleared power errors}
\begin{equation}\label{eq:smallerror}
 d_i=Cm_i-\frac{C}{n_j^b}\,n_j^{b+k_i}.
\end{equation}
This is an integer: $b+k_i\ge0$ and $C/n_j^b$ is an integer. The Pell
congruence \eqref{eq:pellcongruence} gives
$C\delta_i\equiv d_i\pmod {T_j}$. Therefore
\begin{equation}\label{eq:smallfraction}
 \sum_{j=0}^1\frac{\sum_{i\in I_j}d_i^2}{T_j}\in\Z.
\end{equation}
Indeed, multiply \eqref{eq:rawfraction} by $C^2$ and replace each
numerator by a congruent one modulo its own denominator. The difference
is an integer. This argument does not divide by $C$ modulo either
modulus and uses no coprimality between the moduli.

The uniform estimates are
\[
 |C|\le b^{2b},\qquad |C/n_j^b|\le b^b,\qquad
 |d_i|\le b^{2b+1}+b^{3b}\le2b^{3b}.
\]
Writing $r=|I_0|+|I_1|\le15$, we obtain
\begin{equation}\label{eq:budget}
 \sum_i d_i^2\le4r b^{6b}<W^{3b+1}\le W^W.
\end{equation}
The strict inequality uses $4r\le60<W$ and $b^2<W$. The last uses
$3b+1\le b^2+1\le W$, valid for $b\ge3$.
Together with \eqref{eq:Tbound}, this shows that the integer in
\eqref{eq:smallfraction} has absolute value strictly less than one.
It is zero. All its summands have a common sign, so every $d_i$ is zero.

Divide the equality $d_i=0$ by $C$ in $\mathbb Q$. This gives
$m_i=n_j^{k_i}$. If $|n_j|\ge2$ and $k_i<0$, then
$m_i n_j^{-k_i}=1$, impossible for an integer $m_i$. This proves (b)
implies (a), including negative bases and negative Pell parameters.

Conversely, true powers and \eqref{eq:pellcongruence} give
$T_j\mid\delta_i$. Equation~\eqref{eq:squareidentity} then gives
$T_j\mid E_j$. Thus $z=E_0/T_0+E_1/T_1$ is an integer satisfying
\eqref{eq:RP}. This proves (a) implies (b).\hfill$\square$

\begin{remark}[Why both hypotheses matter]
Common sign prevents cancellation: $1/2+1/(-2)=0$ has nonzero
numerators. A strict bound prevents a nonzero integer sum:
$1/2+1/2=1$. These are counterexamples to weakened rational arguments,
not to the theorem. Neither pairwise coprimality nor individual
divisibility can simply be read off a rational sum in the absence of
the size and sign argument.
\end{remark}

\subsection{Replacing the old four rows}

The old same-base certificates use two private coordinates $K_j,L_j$
per base and the rows
\begin{align}
 A_j&=K_j-2H_j-(2A-n_j)L_j,\label{eq:oldA}\\
 B_j&=S_j-n_jK_j+L_j.\label{eq:oldB}
\end{align}
The exact polynomial identity is
\begin{equation}\label{eq:eliminate}
 B_j+n_jA_j=E_j-T_jL_j.
\end{equation}
Old zero witnesses give the new quotient $z=L_0+L_1$. Conversely, the
new theorem gives true powers and hence individual integrality of
$L_j=E_j/T_j$. Restore
$K_j=2H_j+(2A-n_j)L_j$. Both old rows then vanish.

This is an equivalence after existentially quantifying private witnesses.
The inverse quotients are integral \emph{on solutions}; they are not
polynomial functions on the full ambient lattice. The report therefore
does not assert an ambient polynomial automorphism for this semantic
compression.

\section{Instantiation and the bootstrap}

\subsection{The four protected rows}

In the principal 218-variable chart, the last four coordinates are
\[
 x_{214}=K_1,\quad x_{215}=L_1,\quad
 x_{216}=K_0,\quad x_{217}=L_0.
\]
They occur only in rows 151--154. Rows 151,152 are the aggregate and
norm for $q+1$; rows 153,154 are those for $q$. The actual affine ports
are $q=x_{95}-1$ and $A=x_{135}+1$. The protected row transformation
does not alter any of these four rows.

Setting the four private coordinates to zero in
\eqref{eq:eliminate} gives exactly \eqref{eq:Eactual}. The parent has
three requests at $q$ and eleven at $q+1$, hence $r=14$. Its Pell
difference coordinates are exact integral changes of coordinates, and
the norm-row simplifications are identities modulo retained Pell rows.
The inherited pool certificate and the exact projection maps retain
that correspondence. The finite builder does not guess port names or
ignore constant offsets: it extracts and verifies their full affine
polynomials from the actual four rows.

\subsection{Avoiding circular sign reasoning}

To apply Theorem~\ref{thm:pool} one must establish $q\ne0,-1$ before
using either pool. The retained singleton requests and the ceiling
argument supply $|B|\ge2$, $H$ odd and nonzero, $x\ge0$, and
$q=(BH)^x$. Therefore $q=1$ if $x=0$; otherwise $|q|\ge2$.
In either case $q\ne0,-1$, so $q$ and $q+1$ are nonzero.

The theorem now recovers all fourteen requests simultaneously. It is
legitimate even in the possible unit-base case $q=1$. The recovered
request $M=q^x$ and the unchanged equation $Bd=M$ rule out $x=0$:
that would imply $B\mid1$ despite $|B|\ge2$. Thus $x\ge1$ and
$|q|\ge2$. The retained integer ratio $q^t v=q^x$ forces $x\ge t$.
At this stage the parent's rounded-factorial, CRT range, and Horn-flow
arguments resume with exactly their required hypotheses.

For completeness, start with the same finite C222 refutation as in the
parent. Choose a sufficiently large cofinal horizon, construct its true
power outputs and canonical Pell pairs, and use the parent's witness
construction. The new quotient is $L_0+L_1$. The coordinate projection
maps then give a solution of \eqref{eq:companion}. No consistency
assumption is used in this conditional witness construction.

\subsection{Degree and dimension}

Deleting four private variables and adjoining one leaves 215 variables.
Deleting four rows and adjoining one leaves 152 residuals. Each $E_j$
is at most cubic and each $T_j$ quadratic, so $\mathcal R$ is at most
quintic. The coefficient of
\[
 v_{214}^{\,2}v_{95}^{\,8}
\]
in the fully expanded \eqref{eq:companion} is exactly one. Thus degree
ten is exact, not merely an upper bound. The other 151 squared
residuals have degree at most four and cannot cancel that term.

\begin{proof}[Proof of Theorem~\ref{thm:main}]
The five applications of Lemma~\ref{lem:radical}, the affine chart, and
the final affine substitution give explicit integer zero-set bijections
between $P_{225}$ and the raw 218-variable system. The protected
invertible row operations preserve its common zero set, proving the
claim for \eqref{eq:principal}. Theorem~\ref{thm:pool}, the preceding
bootstrap, and the two witness directions after
\eqref{eq:eliminate} prove that \eqref{eq:companion} is solvable exactly
when \eqref{eq:principal} is. Compose with the retained parent arithmetic
theorem and the C222--ZF--ZFC refutation translations. All steps concern
finite syntactic objects, integer arithmetic, recurrences, and finite
constructions, and formalize in ZFC. Negation gives the stated
consistency equivalences.
\end{proof}

\section{What was verified}

\subsection{Independent exact checks}

The principal checker uses only Python's standard library. It imports no
construction or optimization module. Its checks are symbolic integer
coefficient identities, not searches for solutions. They cover:
\begin{itemize}
\item every common-radical condition, complete lift and inverse map,
 and every projected residual;
\item the affine row consequence, the integral hyperplane chart, its
 inverse, and the final unit-coefficient substitution;
\item all 669 residual row operations and their exact inverses;
\item privacy of the four removed coordinates, literal preservation of
 the four pool rows, and every pool elimination/recombination identity;
\item all coefficients of both expanded final polynomials, exact degrees,
 coefficient heights, term counts, and occurrence of every coordinate;
\item deliberate corruptions of maps, source coefficients, moduli,
 row indices, and reindexing data.
\end{itemize}
The separate arithmetic checker verifies the Pell congruence and rational
identities on finite examples with both signs and negative indices. Its
396 signed-congruence cases include 180 negative indices. Eight true-power
and eight perturbed-output cases use the full theorem hypotheses with
$W=258$ and $|A|=258^{258}+1$. The proof of the unbounded theorem is
Section~4, not an extrapolation from these fixtures.

The immediate 225-variable parent's verifier was replayed with a rebuild
of its 20 generated JSON stage files. The restored source closure was
authenticated against its retained hashes and original pinned archives.
The normal and optimized Python checks of the new work agree; acceptance
does not rely on \code{assert}, which optimized Python would remove.
Detailed receipts and the exact verification scope are in
\code{verification/}.

\subsection{Reproducibility and evaluation}

From the extracted archive directory, with Python 3.10 or later:
\begin{verbatim}
python3 -B verify.py
python3 -B -O verify.py
python3 -B verify.py --rebuild
python3 -B code/evaluate.py quartic218 point.json
python3 -B code/evaluate.py decic215 point.json
\end{verbatim}
An evaluator input is exactly one JSON list of 218 or 215 integers,
respectively. It evaluates that specified tuple. It does not decide
whether a zero exists. The verifier works offline and creates no network
requests. The report can be rebuilt by running pdfLaTeX twice in
\code{report/}; the generated residual and metric inputs are included.

\subsection{Proof boundaries}

An exact coefficient audit does not prove the general Pell or logical
calibration theorem. Conversely, a correct general theorem does not
authenticate an accidentally mistyped coefficient list. The archive
provides both the ordinary proofs and the separate finite checks.
The independent reviews are separately assigned AI reviews, not external
human peer review. The transitive Lean/Rocq closure was not freshly
compiled in this iteration, and no theorem of consistency itself was
proved.

\section{Novelty, comparison, and the useful next target}

The old same-base pooling of package 158 uses divisibility of a sum of
squared errors by \emph{one} modulus~\cite{parent158}. The present
replacement uses integrality of a rational sum with \emph{two different}
moduli. Clearing signed indices, proving a common sign, and budgeting
the two denominators against one frame are the extra steps that make
this work. The existing reciprocal-factorial CRT construction uses
moduli $1+F^2/i$ and a binomial product identity~\cite{reciprocal660};
it is a different mechanism from reciprocal pooling of Pell-error
fractions. We found no instance of the new two-base pooling theorem in
the examined snapshot. That is a scoped comparison, not a claim to
have proved historical priority over all mathematical literature.

The seven quartic eliminations combine methods already present in the
repository with a newer source. Their value is the complete, checked
new endpoint and its explicit inverse maps. The affine chart avoids
treating rational substitution as if it automatically preserved integer
solutions. The protected residual basis makes the semantic power-pool
interface stable while changing the surrounding quartic.

The generic principle is reusable: a finite sum of nonnegative rational
errors with denominators of one sign vanishes when it is an integer and
its absolute value is below one. Here the subtlety is that the actual
Pell coordinates are not small; the signed-index clearing and congruence
argument replace them by small errors before applying that principle.
This suggests a concrete further compiler target: merge additional
certificate blocks only after proving a compatible sign orientation and
a common cleared-error budget. The present archive claims the completed
two-base instance, not a completed arbitrary-base compiler.

For a next reduction, the most valuable question is whether the same
cross-base information can be realized with a lower-degree residual
without spending back the three saved variables. A stronger global
``simplest'' claim would require a declared complexity order and a
comparison across the coefficient-heavy and positive-domain branches.
Neither an undecidability theorem nor the absence of further reductions
in one linear-algebra search supplies such a lower bound.

\appendix
\section{Every quadratic residual of the principal equation}
\label{app:residuals}

All variables below are $x_0,\ldots,x_{217}\in\Z$. The polynomial is
the sum of the squares of exactly these 155 rows. Indices are zero-based
both here and in the JSON. In particular the last four rows are the
protected two-base interface used in \eqref{eq:Eactual}.

\begingroup\small
\setlength{\abovedisplayskip}{4pt}\setlength{\belowdisplayskip}{4pt}
% Generated from equations/quartic218/equation.json. Actual math is parsed by check_report.py.
% RESIDUAL 0
\begin{equation*}\begin{aligned}
Q_{0}={}& -x_{2}^{2}+x_{3}
\end{aligned}\end{equation*}
% RESIDUAL 1
\begin{equation*}\begin{aligned}
Q_{1}={}& x_{0}-x_{0}x_{5}-x_{5}+x_{8}
\end{aligned}\end{equation*}
% RESIDUAL 2
\begin{equation*}\begin{aligned}
Q_{2}={}& -x_{2}-x_{5}x_{6}+x_{7}x_{8}
\end{aligned}\end{equation*}
% RESIDUAL 3
\begin{equation*}\begin{aligned}
Q_{3}={}& -1+x_{1}-x_{7}x_{8}
\end{aligned}\end{equation*}
% RESIDUAL 4
\begin{equation*}\begin{aligned}
Q_{4}={}& 1+x_{3}-x_{4}x_{19}+x_{17}
\end{aligned}\end{equation*}
% RESIDUAL 5
\begin{equation*}\begin{aligned}
Q_{5}={}& -1-x_{1}+x_{7}x_{8}+x_{10}-x_{20}^{2}+x_{21}
\end{aligned}\end{equation*}
% RESIDUAL 6
\begin{equation*}\begin{aligned}
Q_{6}={}& -4x_{9}^{2}-4x_{9}x_{10}-8x_{9}x_{20}-16x_{9}x_{21}-2x_{10}-x_{10}x_{20}-8x_{10}x_{21}\\
& -x_{11}-x_{11}^{2}+x_{11}x_{20}+3x_{20}+x_{20}^{2}-16x_{20}x_{21}+5x_{21}+x_{22}+x_{72}
\end{aligned}\end{equation*}
% RESIDUAL 7
\begin{equation*}\begin{aligned}
Q_{7}={}& -2-15x_{0}-5x_{0}x_{5}-x_{1}+x_{3}-x_{3}x_{140}-x_{4}x_{19}-3x_{5}+x_{5}x_{17}\\
& +x_{7}x_{8}+5x_{8}-240x_{9}-56x_{9}^{2}+8x_{9}x_{10}+16x_{9}x_{20}-128x_{9}x_{21}\\
& -64x_{9}x_{22}+87x_{10}+44x_{10}^{2}+42x_{10}x_{11}+71x_{10}x_{20}-32x_{10}x_{21}\\
& -24x_{10}x_{22}+56x_{10}x_{24}-8x_{10}x_{25}+8x_{10}x_{35}-8x_{10}x_{53}+22x_{11}\\
& +14x_{11}^{2}-6x_{11}x_{12}-3x_{11}x_{20}-28x_{11}x_{21}+16x_{11}x_{24}-16x_{11}x_{53}\\
& +2x_{12}-2x_{12}^{2}-50x_{20}-96x_{20}^{2}+32x_{20}x_{21}-48x_{20}x_{22}-8x_{20}x_{25}\\
& +8x_{20}x_{35}-16x_{20}x_{53}-188x_{21}-60x_{21}^{2}-128x_{21}x_{22}+16x_{21}x_{24}\\
& -16x_{21}x_{25}-51x_{22}-4x_{22}^{2}+11x_{23}+16x_{24}-16x_{24}^{2}+19x_{25}+2x_{26}\\
& -8x_{26}^{2}+12x_{27}+10x_{29}-2x_{30}-4x_{30}^{2}+11x_{31}-2x_{32}+10x_{33}+3x_{34}\\
& +20x_{35}-8x_{35}^{2}+12x_{36}+5x_{37}+10x_{39}+3x_{40}+10x_{42}-8x_{53}-8x_{53}^{2}\\
& +2x_{54}-5x_{60}+12x_{61}-8x_{61}^{2}+12x_{62}+9x_{72}-x_{72}^{2}+x_{73}-x_{91}\\
& -x_{135}x_{176}-x_{139}-x_{140}+x_{142}x_{143}-x_{176}+x_{177}+20x_{185}x_{198}\\
& +20x_{188}^{2}-3x_{192}+2x_{192}x_{193}-2x_{200}
\end{aligned}\end{equation*}
% RESIDUAL 8
\begin{equation*}\begin{aligned}
Q_{8}={}& -x_{10}-x_{10}^{2}+2x_{10}x_{11}-2x_{10}x_{20}-x_{11}^{2}+2x_{11}x_{20}-x_{21}+x_{24}
\end{aligned}\end{equation*}
% RESIDUAL 9
\begin{equation*}\begin{aligned}
Q_{9}={}& -1-x_{10}-3x_{11}^{2}+2x_{11}x_{12}+x_{12}^{2}+2x_{20}-x_{21}+x_{24}+x_{25}-x_{40}\\
& +2x_{72}
\end{aligned}\end{equation*}
% RESIDUAL 10
\begin{equation*}\begin{aligned}
Q_{10}={}& -x_{0}-3x_{9}+32x_{9}x_{21}-2x_{10}+11x_{10}^{2}+14x_{10}x_{11}-9x_{10}x_{20}\\
& -8x_{10}x_{21}-4x_{10}x_{24}-24x_{10}x_{25}+16x_{10}x_{44}+8x_{10}x_{53}-14x_{11}\\
& -14x_{11}^{2}+4x_{11}x_{12}-27x_{11}x_{20}-24x_{11}x_{21}-24x_{11}x_{24}-8x_{11}x_{25}\\
& +20x_{11}x_{28}-4x_{11}x_{44}+16x_{11}x_{53}+2x_{12}^{2}-23x_{20}-6x_{20}^{2}\\
& -16x_{20}x_{21}-16x_{20}x_{24}-48x_{20}x_{25}+16x_{20}x_{44}+16x_{20}x_{53}-2x_{21}\\
& -8x_{21}x_{25}-x_{22}+2x_{24}-8x_{24}^{2}-8x_{24}x_{25}-16x_{24}x_{44}-17x_{25}\\
& +8x_{25}^{2}+16x_{25}x_{28}+x_{26}-4x_{26}^{2}+x_{27}+x_{28}+4x_{28}^{2}+x_{34}+x_{37}\\
& -2x_{40}-x_{43}-6x_{44}-4x_{44}^{2}+x_{45}+8x_{53}+8x_{53}^{2}-2x_{54}-2x_{60}-3x_{72}\\
& +x_{188}^{2}
\end{aligned}\end{equation*}
% RESIDUAL 11
\begin{equation*}\begin{aligned}
Q_{11}={}& -x_{0}+x_{10}^{2}+5x_{10}x_{20}+x_{11}+4x_{11}^{2}-2x_{11}x_{12}-x_{11}x_{20}-x_{12}^{2}\\
& +11x_{20}+2x_{20}^{2}-8x_{20}x_{24}+x_{21}+15x_{24}-4x_{24}^{2}-x_{25}-4x_{26}^{2}+x_{27}\\
& +x_{40}-x_{44}-3x_{72}+x_{188}^{2}
\end{aligned}\end{equation*}
% RESIDUAL 12
\begin{equation*}\begin{aligned}
Q_{12}={}& -2x_{0}-3x_{0}^{2}-2x_{0}x_{3}-x_{1}^{2}-x_{2}^{2}-x_{3}-2x_{3}^{2}-x_{4}^{2}-2x_{5}^{2}+x_{10}\\
& +x_{10}^{2}-2x_{10}x_{11}+2x_{10}x_{20}+14x_{11}-5x_{11}^{2}+2x_{11}x_{12}-16x_{11}x_{24}\\
& -16x_{11}x_{25}-x_{13}^{2}-x_{14}^{2}-x_{15}^{2}-x_{16}^{2}-x_{17}^{2}-x_{18}^{2}-4x_{20}+2x_{20}^{2}\\
& +x_{21}+4x_{24}-4x_{24}^{2}-8x_{24}x_{25}+5x_{25}-4x_{25}^{2}+x_{28}-x_{40}-4x_{72}\\
& +x_{95}-2x_{95}^{2}-x_{96}^{2}-x_{97}^{2}-x_{101}^{2}-x_{104}^{2}-x_{107}^{2}-x_{110}^{2}-x_{113}^{2}\\
& -x_{114}^{2}-x_{117}^{2}-x_{120}^{2}-x_{121}^{2}-x_{125}^{2}-x_{128}^{2}-x_{131}^{2}+x_{134}+x_{140}^{2}\\
& -x_{146}x_{147}+x_{149}^{2}-2x_{152}+x_{152}^{2}+x_{155}^{2}+x_{158}^{2}+x_{161}^{2}+x_{164}^{2}\\
& +x_{167}^{2}+x_{170}^{2}+x_{173}^{2}+x_{175}+x_{176}^{2}+x_{176}x_{187}+x_{179}^{2}+x_{182}^{2}\\
& -x_{185}^{2}-x_{187}^{2}+x_{205}^{2}+x_{206}^{2}+x_{207}+x_{207}^{2}
\end{aligned}\end{equation*}
% RESIDUAL 13
\begin{equation*}\begin{aligned}
Q_{13}={}& 2+x_{1}-x_{3}+x_{4}x_{19}-x_{7}x_{8}-x_{9}-2x_{10}-x_{10}^{2}+4x_{10}x_{11}\\
& +4x_{10}x_{40}+23x_{11}-2x_{11}x_{12}+4x_{11}x_{20}-16x_{11}x_{21}+8x_{11}x_{25}\\
& +4x_{11}x_{28}+8x_{11}x_{40}-x_{12}^{2}-x_{17}-4x_{20}+2x_{20}^{2}+8x_{20}x_{25}\\
& +8x_{20}x_{40}+8x_{21}-16x_{21}x_{25}-8x_{21}x_{28}-4x_{24}+8x_{24}x_{25}+4x_{25}\\
& +4x_{25}^{2}-4x_{26}^{2}+x_{27}+12x_{28}+x_{29}-3x_{40}+4x_{40}^{2}-x_{41}-x_{43}-x_{44}\\
& +4x_{72}
\end{aligned}\end{equation*}
% RESIDUAL 14
\begin{equation*}\begin{aligned}
Q_{14}={}& -4x_{0}-x_{0}x_{5}-x_{5}+x_{8}-2x_{9}+4x_{9}^{2}+4x_{9}x_{10}+8x_{9}x_{20}\\
& -16x_{9}x_{21}-9x_{10}-2x_{10}^{2}-4x_{10}x_{11}-4x_{10}x_{20}-8x_{10}x_{24}\\
& +4x_{10}x_{25}+16x_{10}x_{44}+4x_{10}x_{53}+20x_{10}x_{56}-3x_{11}+7x_{11}^{2}\\
& -2x_{11}x_{12}-4x_{11}x_{20}+36x_{11}x_{21}+24x_{11}x_{24}-4x_{11}x_{44}+8x_{11}x_{53}\\
& -x_{12}^{2}-6x_{20}-3x_{20}^{2}+32x_{20}x_{21}-16x_{20}x_{24}+4x_{20}x_{25}+16x_{20}x_{44}\\
& +8x_{20}x_{53}+20x_{20}x_{56}+4x_{21}+8x_{21}x_{25}+9x_{24}-4x_{24}^{2}-16x_{24}x_{44}\\
& -16x_{24}x_{56}-3x_{25}-x_{26}-8x_{26}^{2}+2x_{27}+x_{30}+4x_{30}^{2}-x_{31}-x_{34}\\
& +x_{37}+x_{40}-9x_{44}-4x_{44}^{2}+x_{45}-12x_{49}-4x_{49}^{2}+4x_{53}+4x_{53}^{2}-x_{54}\\
& -4x_{56}^{2}+x_{57}-x_{60}+14x_{72}+5x_{188}^{2}
\end{aligned}\end{equation*}
% RESIDUAL 15
\begin{equation*}\begin{aligned}
Q_{15}={}& 2+2x_{0}+2x_{1}-x_{2}-x_{3}+x_{4}x_{19}-x_{5}x_{6}-x_{7}x_{8}+2x_{9}\\
& -32x_{9}x_{21}+14x_{10}-17x_{10}^{2}-2x_{10}x_{11}+8x_{10}x_{21}+4x_{10}x_{25}\\
& +4x_{10}x_{53}+16x_{10}x_{60}-28x_{11}-11x_{11}^{2}+4x_{11}x_{12}-10x_{11}x_{20}\\
& -40x_{11}x_{24}+8x_{11}x_{25}+4x_{11}x_{28}+8x_{11}x_{53}+4x_{11}x_{58}-2x_{12}+x_{12}^{2}\\
& -x_{17}-6x_{20}+11x_{20}^{2}+16x_{20}x_{21}+4x_{20}x_{25}+8x_{20}x_{53}-8x_{21}x_{24}\\
& -8x_{21}x_{25}-8x_{21}x_{28}+8x_{21}x_{60}+x_{22}-2x_{23}+2x_{24}-12x_{24}^{2}+13x_{25}\\
& -x_{26}-8x_{26}^{2}+12x_{28}-x_{29}+x_{30}-2x_{31}+x_{32}-2x_{33}-x_{34}-4x_{35}\\
& -2x_{36}-x_{37}-2x_{39}-2x_{40}-2x_{42}-x_{43}-2x_{44}+4x_{53}+4x_{53}^{2}-x_{54}\\
& +4x_{58}^{2}-x_{59}-10x_{60}+4x_{60}^{2}-4x_{61}-2x_{62}-18x_{72}+x_{72}^{2}-x_{73}\\
& +x_{135}x_{176}+x_{176}-x_{177}-4x_{185}x_{198}-2x_{188}^{2}+x_{189}^{2}+x_{190}^{2}+x_{192}\\
& -x_{199}^{2}
\end{aligned}\end{equation*}
% RESIDUAL 16
\begin{equation*}\begin{aligned}
Q_{16}={}& 1+x_{0}+x_{1}-x_{3}+x_{4}x_{19}-x_{7}x_{8}+17x_{10}-2x_{10}^{2}-6x_{10}x_{11}\\
& -7x_{10}x_{20}-16x_{10}x_{24}+4x_{11}-4x_{11}^{2}+2x_{11}x_{12}+3x_{11}x_{20}\\
& -4x_{11}x_{21}+8x_{11}x_{24}+16x_{11}x_{25}+x_{12}^{2}-x_{17}-6x_{20}+13x_{20}^{2}\\
& +8x_{20}x_{24}-4x_{21}-8x_{21}x_{24}-2x_{23}+5x_{24}+4x_{24}^{2}+8x_{24}x_{25}-4x_{25}\\
& +4x_{25}^{2}+4x_{26}^{2}-3x_{27}-x_{28}-2x_{29}-2x_{31}+x_{32}-2x_{33}-4x_{35}-2x_{36}\\
& -2x_{39}-x_{40}-2x_{42}+x_{44}-x_{60}-4x_{61}-2x_{62}+9x_{72}+x_{135}x_{176}+x_{176}\\
& -x_{177}-4x_{185}x_{198}-x_{188}^{2}+x_{189}^{2}+x_{190}^{2}+x_{192}-x_{199}^{2}
\end{aligned}\end{equation*}
% RESIDUAL 17
\begin{equation*}\begin{aligned}
Q_{17}={}& 1+4x_{0}+3x_{0}^{2}+2x_{0}x_{3}+x_{0}x_{5}+x_{1}+x_{1}^{2}+x_{3}+2x_{3}^{2}+x_{4}^{2}\\
& +x_{4}x_{19}+x_{5}+2x_{5}^{2}-x_{7}x_{8}-x_{8}-2x_{9}-6x_{10}+2x_{10}^{2}-28x_{10}x_{11}\\
& -29x_{10}x_{20}-8x_{10}x_{24}+8x_{10}x_{25}-32x_{10}x_{32}-4x_{10}x_{35}-4x_{10}x_{40}\\
& +16x_{10}x_{44}+4x_{10}x_{53}-18x_{11}+9x_{11}^{2}-2x_{11}x_{12}+13x_{11}x_{20}\\
& +4x_{11}x_{21}-32x_{11}x_{24}+8x_{11}x_{25}-40x_{11}x_{28}-40x_{11}x_{32}-8x_{11}x_{40}\\
& -4x_{11}x_{44}+8x_{11}x_{53}+x_{13}^{2}+x_{14}^{2}+x_{15}^{2}+x_{16}^{2}-x_{17}+x_{17}^{2}+x_{18}^{2}\\
& -7x_{20}+5x_{20}^{2}-40x_{20}x_{24}+16x_{20}x_{25}-64x_{20}x_{32}-4x_{20}x_{35}\\
& -8x_{20}x_{40}+16x_{20}x_{44}+8x_{20}x_{53}+12x_{21}-8x_{21}x_{24}-13x_{24}-32x_{24}^{2}\\
& +8x_{24}x_{25}-16x_{24}x_{32}-16x_{24}x_{44}+11x_{25}-16x_{25}^{2}-32x_{25}x_{28}\\
& -8x_{25}x_{32}+12x_{26}^{2}-3x_{27}-4x_{28}-8x_{28}^{2}-2x_{30}-4x_{30}^{2}+x_{31}-3x_{32}\\
& -4x_{32}^{2}+x_{33}+4x_{35}^{2}-x_{36}+x_{37}+5x_{40}-4x_{40}^{2}+x_{41}+2x_{43}+8x_{44}\\
& -4x_{44}^{2}+x_{45}+4x_{53}+4x_{53}^{2}-x_{54}-3x_{60}+35x_{72}-x_{95}+2x_{95}^{2}+x_{96}^{2}\\
& +x_{97}^{2}+x_{101}^{2}+x_{104}^{2}+x_{107}^{2}+x_{110}^{2}+x_{113}^{2}+x_{114}^{2}+x_{117}^{2}+x_{120}^{2}\\
& +x_{121}^{2}+x_{125}^{2}+x_{128}^{2}+x_{131}^{2}-x_{134}-x_{140}^{2}+x_{146}x_{147}-x_{149}^{2}\\
& +2x_{152}-x_{152}^{2}-x_{155}^{2}-x_{158}^{2}-x_{161}^{2}-x_{164}^{2}-x_{167}^{2}-x_{170}^{2}-x_{173}^{2}\\
& -x_{175}-x_{176}^{2}-x_{176}x_{187}-x_{179}^{2}-x_{182}^{2}+x_{185}^{2}+x_{187}^{2}-3x_{188}^{2}\\
& -x_{205}^{2}-x_{206}^{2}-x_{207}-x_{207}^{2}
\end{aligned}\end{equation*}
% RESIDUAL 18
\begin{equation*}\begin{aligned}
Q_{18}={}& 1+x_{9}-4x_{21}^{2}
\end{aligned}\end{equation*}
% RESIDUAL 19
\begin{equation*}\begin{aligned}
Q_{19}={}& -2-x_{1}+x_{7}x_{8}-8x_{9}^{2}-8x_{9}x_{10}-16x_{9}x_{20}+x_{10}^{2}-2x_{10}x_{11}\\
& -2x_{11}+3x_{11}^{2}+4x_{20}-2x_{21}+x_{22}-2x_{24}-x_{25}+x_{34}
\end{aligned}\end{equation*}
% RESIDUAL 20
\begin{equation*}\begin{aligned}
Q_{20}={}& 2+12x_{0}-3x_{0}^{2}-2x_{0}x_{3}+5x_{0}x_{5}+x_{1}-x_{1}^{2}-x_{2}^{2}-x_{3}-2x_{3}^{2}\\
& +x_{3}x_{140}-x_{4}^{2}+3x_{5}-2x_{5}^{2}-x_{5}x_{17}-x_{7}x_{8}-5x_{8}-190x_{9}\\
& -56x_{9}^{2}+8x_{9}x_{10}+16x_{9}x_{20}+64x_{9}x_{21}-64x_{9}x_{34}-68x_{10}-86x_{10}^{2}\\
& -32x_{10}x_{11}-27x_{10}x_{20}+104x_{10}x_{21}-36x_{10}x_{24}+24x_{10}x_{25}\\
& +32x_{10}x_{32}-24x_{10}x_{34}-4x_{10}x_{35}-4x_{10}x_{40}-4x_{10}x_{45}-4x_{10}x_{48}\\
& +8x_{10}x_{53}+16x_{10}x_{60}-15x_{11}-45x_{11}^{2}+18x_{11}x_{12}-5x_{11}x_{20}\\
& +28x_{11}x_{21}-16x_{11}x_{24}+24x_{11}x_{25}+16x_{11}x_{28}+40x_{11}x_{32}\\
& -8x_{11}x_{40}+8x_{11}x_{45}+8x_{11}x_{48}+16x_{11}x_{53}+4x_{11}x_{58}-4x_{12}\\
& +6x_{12}^{2}-x_{13}^{2}-x_{14}^{2}-x_{15}^{2}-x_{16}^{2}+x_{17}-x_{17}^{2}-x_{18}^{2}+36x_{20}\\
& +80x_{20}^{2}-48x_{20}x_{21}+24x_{20}x_{24}+24x_{20}x_{25}+64x_{20}x_{32}-48x_{20}x_{34}\\
& -4x_{20}x_{35}-8x_{20}x_{40}-4x_{20}x_{45}-4x_{20}x_{48}+16x_{20}x_{53}+150x_{21}\\
& -48x_{21}^{2}-16x_{21}x_{24}+40x_{21}x_{25}+8x_{21}x_{28}+128x_{21}x_{34}+8x_{21}x_{60}\\
& +5x_{22}-10x_{23}-28x_{24}+36x_{24}^{2}+16x_{24}x_{32}+8x_{25}+8x_{25}^{2}\\
& +16x_{25}x_{28}+8x_{25}x_{32}-3x_{26}-16x_{26}^{2}-6x_{27}-11x_{28}+4x_{28}^{2}-11x_{29}\\
& +4x_{30}+4x_{30}^{2}-11x_{31}+6x_{32}+4x_{32}^{2}-11x_{33}-53x_{34}-4x_{34}^{2}-19x_{35}\\
& +4x_{35}^{2}-11x_{36}-13x_{37}-10x_{39}-5x_{40}-4x_{40}^{2}+x_{41}-10x_{42}-16x_{44}\\
& +4x_{45}^{2}-x_{46}+4x_{48}^{2}-x_{49}+8x_{53}+8x_{53}^{2}-2x_{54}+4x_{58}^{2}-x_{59}\\
& +5x_{60}+4x_{60}^{2}-12x_{61}+8x_{61}^{2}-12x_{62}-77x_{72}+2x_{72}^{2}-2x_{73}+x_{91}\\
& +2x_{95}-2x_{95}^{2}-x_{96}^{2}-x_{97}^{2}-x_{101}^{2}-x_{104}^{2}-x_{107}^{2}-x_{110}^{2}-x_{113}^{2}\\
& -x_{114}^{2}-x_{117}^{2}-x_{120}^{2}-x_{121}^{2}-x_{125}^{2}-x_{128}^{2}-x_{131}^{2}+x_{134}-x_{135}\\
& +x_{135}x_{152}+x_{139}+x_{140}+x_{140}^{2}-x_{142}x_{143}+x_{149}^{2}-x_{152}+x_{152}^{2}\\
& -x_{153}+x_{155}^{2}+x_{158}^{2}+x_{161}^{2}+x_{164}^{2}+x_{167}^{2}+x_{170}^{2}+x_{173}^{2}+x_{176}^{2}\\
& +x_{179}^{2}+x_{182}^{2}-x_{185}^{2}-20x_{185}x_{198}-x_{187}^{2}-19x_{188}^{2}+x_{189}^{2}+x_{190}^{2}\\
& +3x_{192}-2x_{192}x_{193}-x_{199}^{2}+2x_{200}+x_{205}^{2}+x_{206}^{2}+x_{207}+x_{207}^{2}
\end{aligned}\end{equation*}
% RESIDUAL 21
\begin{equation*}\begin{aligned}
Q_{21}={}& -x_{10}^{2}-2x_{10}x_{20}+4x_{10}x_{35}-x_{20}^{2}+4x_{20}x_{35}-4x_{35}^{2}+x_{36}
\end{aligned}\end{equation*}
% RESIDUAL 22
\begin{equation*}\begin{aligned}
Q_{22}={}& 1-2x_{9}+15x_{10}^{2}+14x_{10}x_{11}-4x_{10}x_{20}-4x_{11}+x_{11}^{2}-4x_{11}x_{20}\\
& +8x_{11}x_{21}+4x_{20}+2x_{20}^{2}+4x_{21}+x_{25}+x_{37}-2x_{60}
\end{aligned}\end{equation*}
% RESIDUAL 23
\begin{equation*}\begin{aligned}
Q_{23}={}& 2x_{0}-14x_{10}-3x_{10}^{2}+2x_{10}x_{11}+4x_{10}x_{20}+8x_{10}x_{24}+32x_{10}x_{37}\\
& -16x_{10}x_{44}+4x_{11}+5x_{11}^{2}-2x_{11}x_{12}+4x_{11}x_{20}-16x_{11}x_{24}\\
& -24x_{11}x_{37}+4x_{11}x_{44}-x_{12}^{2}+6x_{20}-5x_{20}^{2}-8x_{20}x_{24}+32x_{20}x_{37}\\
& -16x_{20}x_{44}+8x_{21}-x_{24}-8x_{24}x_{37}+16x_{24}x_{44}-8x_{26}^{2}+2x_{27}+4x_{37}\\
& -4x_{37}^{2}+x_{38}+x_{40}-x_{44}+4x_{44}^{2}-x_{45}+2x_{72}-2x_{188}^{2}
\end{aligned}\end{equation*}
% RESIDUAL 24
\begin{equation*}\begin{aligned}
Q_{24}={}& 7x_{0}+3x_{0}^{2}+2x_{0}x_{3}+2x_{0}x_{5}+x_{1}^{2}+x_{2}^{2}+x_{3}+2x_{3}^{2}+x_{4}^{2}\\
& +2x_{5}+2x_{5}^{2}-2x_{8}+9x_{10}+7x_{10}^{2}-12x_{10}x_{11}-6x_{10}x_{20}-4x_{10}x_{24}\\
& -12x_{10}x_{25}-32x_{10}x_{38}-16x_{10}x_{44}+4x_{10}x_{45}-21x_{11}+9x_{11}^{2}\\
& -4x_{11}x_{12}+12x_{11}x_{20}-4x_{11}x_{21}-32x_{11}x_{24}+8x_{11}x_{25}-24x_{11}x_{38}\\
& +4x_{11}x_{44}-8x_{11}x_{45}-x_{12}^{2}+x_{13}^{2}+x_{14}^{2}+x_{15}^{2}+x_{16}^{2}+x_{17}^{2}\\
& +x_{18}^{2}+12x_{20}-20x_{20}x_{25}-64x_{20}x_{38}-16x_{20}x_{44}+4x_{20}x_{45}-10x_{21}\\
& -8x_{21}x_{24}+7x_{24}-8x_{24}^{2}+16x_{24}x_{44}-3x_{25}-8x_{25}x_{38}-4x_{26}^{2}+x_{27}\\
& -x_{30}-4x_{30}^{2}+x_{31}+x_{32}+4x_{38}-4x_{38}^{2}+x_{39}+2x_{40}+2x_{44}+4x_{44}^{2}\\
& -x_{45}-4x_{45}^{2}+x_{46}-x_{60}+6x_{72}-2x_{95}+2x_{95}^{2}+x_{96}^{2}+x_{97}^{2}+x_{101}^{2}\\
& +x_{104}^{2}+x_{107}^{2}+x_{110}^{2}+x_{113}^{2}+x_{114}^{2}+x_{117}^{2}+x_{120}^{2}+x_{121}^{2}+x_{125}^{2}\\
& +x_{128}^{2}+x_{131}^{2}-x_{134}+x_{135}-x_{135}x_{152}-x_{140}^{2}-x_{149}^{2}+x_{152}-x_{152}^{2}\\
& +x_{153}-x_{155}^{2}-x_{158}^{2}-x_{161}^{2}-x_{164}^{2}-x_{167}^{2}-x_{170}^{2}-x_{173}^{2}-x_{176}^{2}\\
& -x_{179}^{2}-x_{182}^{2}+x_{185}^{2}+x_{187}^{2}-7x_{188}^{2}-x_{205}^{2}-x_{206}^{2}-x_{207}-x_{207}^{2}
\end{aligned}\end{equation*}
% RESIDUAL 25
\begin{equation*}\begin{aligned}
Q_{25}={}& -x_{11}^{2}-2x_{11}x_{12}-x_{12}^{2}+x_{40}
\end{aligned}\end{equation*}
% RESIDUAL 26
\begin{equation*}\begin{aligned}
Q_{26}={}& 1-x_{0}+x_{1}-x_{7}x_{8}+x_{10}^{2}+2x_{10}x_{20}-4x_{10}x_{40}-2x_{11}+4x_{11}^{2}\\
& +2x_{11}x_{20}-8x_{11}x_{40}+2x_{20}-8x_{20}x_{25}-8x_{20}x_{40}+2x_{21}+6x_{24}\\
& -8x_{24}x_{25}+20x_{25}-4x_{25}^{2}+4x_{26}^{2}-x_{27}+4x_{40}-4x_{40}^{2}+x_{41}+x_{44}\\
& -2x_{72}+x_{188}^{2}
\end{aligned}\end{equation*}
% RESIDUAL 27
\begin{equation*}\begin{aligned}
Q_{27}={}& -2+2x_{0}-x_{1}+x_{7}x_{8}-32x_{9}x_{21}-7x_{10}+2x_{10}x_{20}+8x_{10}x_{21}\\
& +12x_{10}x_{24}-8x_{10}x_{25}-28x_{10}x_{40}-16x_{10}x_{41}-8x_{10}x_{53}+25x_{11}\\
& +6x_{11}^{2}-6x_{11}x_{12}+10x_{11}x_{20}+16x_{11}x_{21}+8x_{11}x_{25}+4x_{11}x_{28}\\
& -24x_{11}x_{40}-16x_{11}x_{41}-16x_{11}x_{53}+2x_{12}-2x_{12}^{2}+14x_{20}+5x_{20}^{2}\\
& +16x_{20}x_{21}+8x_{20}x_{24}+8x_{20}x_{25}-56x_{20}x_{40}-32x_{20}x_{41}-16x_{20}x_{53}\\
& +10x_{21}-8x_{21}x_{25}-8x_{21}x_{28}+x_{22}-3x_{24}-8x_{24}^{2}+24x_{24}x_{25}\\
& -25x_{25}-32x_{25}x_{40}-16x_{25}x_{41}-x_{26}-16x_{26}^{2}+4x_{27}+12x_{28}+x_{29}\\
& -x_{34}-5x_{40}-12x_{40}^{2}-16x_{40}x_{41}-3x_{41}-4x_{41}^{2}+x_{42}-x_{43}-2x_{44}\\
& -8x_{53}-8x_{53}^{2}+2x_{54}+12x_{72}-x_{72}^{2}+x_{73}-2x_{188}^{2}
\end{aligned}\end{equation*}
% RESIDUAL 28
\begin{equation*}\begin{aligned}
Q_{28}={}& -1+x_{0}-x_{10}-x_{10}^{2}-2x_{10}x_{11}-4x_{10}x_{20}-13x_{11}-5x_{11}^{2}+2x_{11}x_{12}\\
& -4x_{11}x_{20}+16x_{11}x_{24}-24x_{11}x_{25}-20x_{11}x_{28}+x_{12}^{2}+2x_{20}-3x_{21}\\
& -9x_{24}+4x_{24}^{2}+8x_{24}x_{25}-8x_{25}-12x_{25}^{2}-16x_{25}x_{28}-x_{28}-4x_{28}^{2}\\
& -x_{40}+x_{43}-2x_{72}-x_{188}^{2}-x_{189}^{2}-x_{190}^{2}+x_{199}^{2}
\end{aligned}\end{equation*}
% RESIDUAL 29
\begin{equation*}\begin{aligned}
Q_{29}={}& -1-7x_{10}-x_{10}^{2}-4x_{10}x_{11}-x_{10}x_{20}+8x_{10}x_{24}-x_{11}+x_{11}^{2}\\
& -3x_{11}x_{20}-9x_{20}-2x_{20}^{2}+8x_{20}x_{24}+2x_{21}+6x_{24}-4x_{24}^{2}+x_{44}\\
& +x_{72}
\end{aligned}\end{equation*}
% RESIDUAL 30
\begin{equation*}\begin{aligned}
Q_{30}={}& -1+7x_{10}-2x_{10}^{2}+2x_{10}x_{11}+2x_{10}x_{20}+16x_{10}x_{24}+16x_{10}x_{44}\\
& -4x_{10}x_{53}-8x_{11}+3x_{11}^{2}-2x_{11}x_{12}+2x_{11}x_{20}-8x_{11}x_{24}-4x_{11}x_{44}\\
& -8x_{11}x_{53}-x_{12}^{2}+16x_{20}+6x_{20}^{2}+8x_{20}x_{24}+16x_{20}x_{44}-8x_{20}x_{53}\\
& -2x_{21}-4x_{24}-12x_{24}^{2}-16x_{24}x_{44}-3x_{25}-4x_{26}^{2}+x_{27}+x_{40}-3x_{44}\\
& -4x_{44}^{2}+x_{45}-4x_{53}-4x_{53}^{2}+x_{54}-2x_{72}
\end{aligned}\end{equation*}
% RESIDUAL 31
\begin{equation*}\begin{aligned}
Q_{31}={}& 1+x_{10}+x_{10}^{2}+2x_{10}x_{20}+4x_{10}x_{45}-2x_{11}^{2}-8x_{11}x_{45}-x_{20}-x_{20}^{2}\\
& +4x_{20}x_{45}+2x_{21}-2x_{24}-4x_{45}^{2}+x_{46}
\end{aligned}\end{equation*}
% RESIDUAL 32
\begin{equation*}\begin{aligned}
Q_{32}={}& -1-4x_{46}-4x_{46}^{2}+x_{47}
\end{aligned}\end{equation*}
% RESIDUAL 33
\begin{equation*}\begin{aligned}
Q_{33}={}& 1-x_{11}^{2}-8x_{11}x_{46}-4x_{11}x_{47}+8x_{46}-8x_{46}^{2}-16x_{46}x_{47}-2x_{47}\\
& -4x_{47}^{2}+x_{48}
\end{aligned}\end{equation*}
% RESIDUAL 34
\begin{equation*}\begin{aligned}
Q_{34}={}& 1+x_{1}-x_{7}x_{8}+x_{10}+x_{10}^{2}+2x_{10}x_{20}+4x_{10}x_{48}-2x_{11}^{2}\\
& -8x_{11}x_{48}-x_{20}-x_{20}^{2}+4x_{20}x_{48}+2x_{21}-2x_{24}-4x_{48}^{2}+x_{49}
\end{aligned}\end{equation*}
% RESIDUAL 35
\begin{equation*}\begin{aligned}
Q_{35}={}& -4x_{43}^{2}-8x_{43}x_{49}-4x_{49}^{2}+x_{50}
\end{aligned}\end{equation*}
% RESIDUAL 36
\begin{equation*}\begin{aligned}
Q_{36}={}& -1-x_{1}+x_{7}x_{8}+3x_{10}+x_{10}^{2}+2x_{10}x_{20}-8x_{10}x_{43}-4x_{10}x_{50}\\
& -2x_{11}+2x_{11}^{2}+2x_{11}x_{20}+10x_{20}+3x_{20}^{2}-16x_{20}x_{43}-8x_{20}x_{50}\\
& -2x_{21}-x_{24}-x_{25}+24x_{43}-12x_{43}^{2}+8x_{43}x_{49}-16x_{43}x_{50}+4x_{49}^{2}\\
& +11x_{50}-4x_{50}^{2}+x_{51}
\end{aligned}\end{equation*}
% RESIDUAL 37
\begin{equation*}\begin{aligned}
Q_{37}={}& -1-4x_{51}-4x_{51}^{2}+x_{52}
\end{aligned}\end{equation*}
% RESIDUAL 38
\begin{equation*}\begin{aligned}
Q_{38}={}& 1+x_{1}-x_{7}x_{8}-x_{11}^{2}-8x_{11}x_{51}-4x_{11}x_{52}+8x_{51}-8x_{51}^{2}\\
& -16x_{51}x_{52}-2x_{52}-4x_{52}^{2}+x_{53}
\end{aligned}\end{equation*}
% RESIDUAL 39
\begin{equation*}\begin{aligned}
Q_{39}={}& 1-2x_{10}-x_{10}^{2}-2x_{10}x_{11}-3x_{10}x_{20}-4x_{10}x_{53}-5x_{11}-2x_{11}^{2}\\
& -5x_{11}x_{20}-8x_{11}x_{53}-5x_{20}-2x_{20}^{2}-8x_{20}x_{53}-x_{25}-4x_{53}-4x_{53}^{2}\\
& +x_{54}-x_{72}
\end{aligned}\end{equation*}
% RESIDUAL 40
\begin{equation*}\begin{aligned}
Q_{40}={}& 2-x_{0}+x_{1}-x_{3}+x_{4}x_{19}-x_{7}x_{8}+x_{10}-x_{17}-x_{20}^{2}+x_{21}-12x_{54}\\
& -4x_{54}^{2}+x_{55}+x_{188}^{2}
\end{aligned}\end{equation*}
% RESIDUAL 41
\begin{equation*}\begin{aligned}
Q_{41}={}& 1+x_{1}-x_{7}x_{8}+4x_{10}+16x_{10}x_{54}+8x_{10}x_{55}+4x_{20}+16x_{20}x_{54}\\
& +8x_{20}x_{55}-4x_{24}-16x_{24}x_{54}-8x_{24}x_{55}-x_{44}+8x_{54}-16x_{54}^{2}\\
& -16x_{54}x_{55}+4x_{55}-4x_{55}^{2}+x_{56}
\end{aligned}\end{equation*}
% RESIDUAL 42
\begin{equation*}\begin{aligned}
Q_{42}={}& 1-x_{0}+x_{1}-x_{7}x_{8}+3x_{10}-3x_{10}^{2}-2x_{10}x_{11}-x_{10}x_{20}-16x_{10}x_{44}\\
& +20x_{10}x_{56}-5x_{11}-3x_{11}^{2}+3x_{11}x_{20}+8x_{11}x_{24}+4x_{11}x_{44}+2x_{20}^{2}\\
& -16x_{20}x_{44}+20x_{20}x_{56}+x_{21}-6x_{24}+4x_{24}^{2}+16x_{24}x_{44}-16x_{24}x_{56}\\
& -x_{25}-x_{44}+4x_{44}^{2}-x_{45}-12x_{49}-4x_{49}^{2}-4x_{56}^{2}+x_{57}+3x_{72}+x_{188}^{2}\\
& +x_{189}^{2}+x_{190}^{2}-x_{199}^{2}
\end{aligned}\end{equation*}
% RESIDUAL 43
\begin{equation*}\begin{aligned}
Q_{43}={}& 2+x_{0}+x_{1}-x_{7}x_{8}+x_{10}-x_{20}^{2}+x_{21}+24x_{49}-8x_{49}^{2}-16x_{49}x_{57}\\
& -2x_{57}-4x_{57}^{2}+x_{58}-x_{188}^{2}
\end{aligned}\end{equation*}
% RESIDUAL 44
\begin{equation*}\begin{aligned}
Q_{44}={}& -1+x_{0}+7x_{10}+6x_{10}x_{11}-x_{10}x_{20}-16x_{10}x_{60}-x_{11}+2x_{11}^{2}\\
& -3x_{11}x_{20}+4x_{11}x_{21}-4x_{11}x_{58}+x_{20}+9x_{20}^{2}-2x_{21}-8x_{21}x_{60}\\
& -x_{24}-4x_{58}^{2}+x_{59}+11x_{60}-4x_{60}^{2}-x_{72}-x_{188}^{2}
\end{aligned}\end{equation*}
% RESIDUAL 45
\begin{equation*}\begin{aligned}
Q_{45}={}& 2+2x_{1}-x_{2}-x_{5}x_{6}-x_{7}x_{8}-3x_{9}+5x_{10}+6x_{10}x_{11}-x_{10}x_{20}\\
& -16x_{10}x_{21}-2x_{11}-x_{11}^{2}-3x_{11}x_{20}+4x_{11}x_{21}+3x_{20}+6x_{20}^{2}+6x_{21}\\
& +x_{25}+x_{37}-x_{60}+x_{72}
\end{aligned}\end{equation*}
% RESIDUAL 46
\begin{equation*}\begin{aligned}
Q_{46}={}& -1+x_{0}-x_{1}+x_{7}x_{8}+2x_{9}+12x_{10}-14x_{10}^{2}+10x_{10}x_{11}+3x_{10}x_{20}\\
& -32x_{10}x_{59}-7x_{11}-7x_{11}x_{20}+8x_{11}x_{21}+8x_{11}x_{58}+5x_{20}+8x_{20}^{2}\\
& -7x_{21}-16x_{21}x_{59}-x_{24}+2x_{25}-x_{37}+8x_{58}^{2}+14x_{59}-4x_{59}^{2}-2x_{60}\\
& +x_{61}+3x_{72}-x_{188}^{2}
\end{aligned}\end{equation*}
% RESIDUAL 47
\begin{equation*}\begin{aligned}
Q_{47}={}& -1-x_{1}+x_{7}x_{8}-4x_{61}-4x_{61}^{2}+x_{62}
\end{aligned}\end{equation*}
% RESIDUAL 48
\begin{equation*}\begin{aligned}
Q_{48}={}& -x_{2}x_{14}+x_{63}
\end{aligned}\end{equation*}
% RESIDUAL 49
\begin{equation*}\begin{aligned}
Q_{49}={}& x_{13}-x_{14}-2x_{23}x_{63}-x_{64}
\end{aligned}\end{equation*}
% RESIDUAL 50
\begin{equation*}\begin{aligned}
Q_{50}={}& x_{13}-x_{14}-2x_{27}x_{63}-x_{65}
\end{aligned}\end{equation*}
% RESIDUAL 51
\begin{equation*}\begin{aligned}
Q_{51}={}& -x_{2}x_{14}-2x_{29}x_{63}+2x_{42}x_{63}-x_{63}-x_{66}+x_{71}
\end{aligned}\end{equation*}
% RESIDUAL 52
\begin{equation*}\begin{aligned}
Q_{52}={}& x_{13}-x_{14}-2x_{31}x_{63}-x_{67}
\end{aligned}\end{equation*}
% RESIDUAL 53
\begin{equation*}\begin{aligned}
Q_{53}={}& x_{13}-x_{14}+2x_{23}x_{63}-2x_{29}x_{63}-2x_{33}x_{63}+x_{64}-x_{66}-x_{68}
\end{aligned}\end{equation*}
% RESIDUAL 54
\begin{equation*}\begin{aligned}
Q_{54}={}& x_{13}-x_{14}-4x_{35}x_{63}-2x_{36}x_{63}-x_{69}
\end{aligned}\end{equation*}
% RESIDUAL 55
\begin{equation*}\begin{aligned}
Q_{55}={}& x_{13}-x_{14}-2x_{39}x_{63}-x_{70}
\end{aligned}\end{equation*}
% RESIDUAL 56
\begin{equation*}\begin{aligned}
Q_{56}={}& -2x_{2}x_{14}+x_{13}-x_{14}-2x_{42}x_{63}-2x_{63}-x_{71}
\end{aligned}\end{equation*}
% RESIDUAL 57
\begin{equation*}\begin{aligned}
Q_{57}={}& -x_{10}+x_{10}^{2}+3x_{10}x_{20}-x_{11}-x_{11}^{2}+x_{11}x_{20}-x_{20}+2x_{20}^{2}+x_{72}
\end{aligned}\end{equation*}
% RESIDUAL 58
\begin{equation*}\begin{aligned}
Q_{58}={}& 1+2x_{12}+x_{12}^{2}-x_{72}^{2}+x_{73}
\end{aligned}\end{equation*}
% RESIDUAL 59
\begin{equation*}\begin{aligned}
Q_{59}={}& -x_{13}^{2}+2x_{13}x_{14}+2x_{13}x_{75}-x_{14}^{2}-2x_{14}x_{75}-x_{73}^{2}+x_{74}-x_{75}^{2}
\end{aligned}\end{equation*}
% RESIDUAL 60
\begin{equation*}\begin{aligned}
Q_{60}={}& x_{2}x_{14}-4x_{61}x_{63}-2x_{62}x_{63}+x_{63}+x_{75}
\end{aligned}\end{equation*}
% RESIDUAL 61
\begin{equation*}\begin{aligned}
Q_{61}={}& x_{10}x_{63}-x_{14}+x_{20}x_{63}+x_{76}x_{77}+x_{78}
\end{aligned}\end{equation*}
% RESIDUAL 62
\begin{equation*}\begin{aligned}
Q_{62}={}& -1+3x_{2}-4x_{2}x_{10}-2x_{2}x_{21}-x_{5}x_{6}+x_{7}x_{8}+x_{76}
\end{aligned}\end{equation*}
% RESIDUAL 63
\begin{equation*}\begin{aligned}
Q_{63}={}& -x_{10}x_{213}-x_{13}-2x_{20}x_{213}+x_{77}+x_{213}
\end{aligned}\end{equation*}
% RESIDUAL 64
\begin{equation*}\begin{aligned}
Q_{64}={}& x_{2}x_{14}-x_{10}x_{213}+x_{13}-x_{14}-x_{20}x_{213}-2x_{21}x_{63}+x_{63}+x_{79}
\end{aligned}\end{equation*}
% RESIDUAL 65
\begin{equation*}\begin{aligned}
Q_{65}={}& -x_{13}x_{64}+x_{14}x_{64}+x_{80}
\end{aligned}\end{equation*}
% RESIDUAL 66
\begin{equation*}\begin{aligned}
Q_{66}={}& -x_{65}x_{80}+x_{81}
\end{aligned}\end{equation*}
% RESIDUAL 67
\begin{equation*}\begin{aligned}
Q_{67}={}& -x_{66}x_{81}+x_{82}
\end{aligned}\end{equation*}
% RESIDUAL 68
\begin{equation*}\begin{aligned}
Q_{68}={}& -x_{67}x_{82}+x_{83}
\end{aligned}\end{equation*}
% RESIDUAL 69
\begin{equation*}\begin{aligned}
Q_{69}={}& -x_{68}x_{83}+x_{84}
\end{aligned}\end{equation*}
% RESIDUAL 70
\begin{equation*}\begin{aligned}
Q_{70}={}& -x_{69}x_{84}+x_{85}
\end{aligned}\end{equation*}
% RESIDUAL 71
\begin{equation*}\begin{aligned}
Q_{71}={}& -x_{70}x_{85}+x_{86}
\end{aligned}\end{equation*}
% RESIDUAL 72
\begin{equation*}\begin{aligned}
Q_{72}={}& -x_{71}x_{86}+x_{87}
\end{aligned}\end{equation*}
% RESIDUAL 73
\begin{equation*}\begin{aligned}
Q_{73}={}& -x_{78}x_{87}+x_{88}
\end{aligned}\end{equation*}
% RESIDUAL 74
\begin{equation*}\begin{aligned}
Q_{74}={}& -x_{79}x_{88}+x_{89}
\end{aligned}\end{equation*}
% RESIDUAL 75
\begin{equation*}\begin{aligned}
Q_{75}={}& -x_{89}^{2}+x_{90}
\end{aligned}\end{equation*}
% RESIDUAL 76
\begin{equation*}\begin{aligned}
Q_{76}={}& -x_{5}-x_{5}x_{17}+x_{17}+x_{91}
\end{aligned}\end{equation*}
% RESIDUAL 77
\begin{equation*}\begin{aligned}
Q_{77}={}& -x_{15}x_{91}+x_{92}
\end{aligned}\end{equation*}
% RESIDUAL 78
\begin{equation*}\begin{aligned}
Q_{78}={}& -2x_{1}^{2}+2x_{1}x_{13}+2x_{1}x_{14}+2x_{1}x_{92}+2x_{1}x_{93}-x_{13}^{2}-2x_{13}x_{92}\\
& -x_{14}^{2}-2x_{14}x_{93}+x_{18}x_{94}-x_{74}x_{90}-x_{92}^{2}-x_{93}^{2}
\end{aligned}\end{equation*}
% RESIDUAL 79
\begin{equation*}\begin{aligned}
Q_{79}={}& -x_{5}x_{16}+x_{16}-x_{16}x_{91}+x_{93}
\end{aligned}\end{equation*}
% RESIDUAL 80
\begin{equation*}\begin{aligned}
Q_{80}={}& -x_{95}x_{98}+x_{98}+x_{100}
\end{aligned}\end{equation*}
% RESIDUAL 81
\begin{equation*}\begin{aligned}
Q_{81}={}& -x_{4}x_{121}+x_{96}x_{203}-x_{99}-x_{100}x_{121}
\end{aligned}\end{equation*}
% RESIDUAL 82
\begin{equation*}\begin{aligned}
Q_{82}={}& x_{101}x_{203}-x_{102}x_{124}-x_{103}
\end{aligned}\end{equation*}
% RESIDUAL 83
\begin{equation*}\begin{aligned}
Q_{83}={}& x_{104}x_{203}-x_{105}x_{124}-x_{106}
\end{aligned}\end{equation*}
% RESIDUAL 84
\begin{equation*}\begin{aligned}
Q_{84}={}& x_{107}x_{203}-x_{108}x_{124}-x_{109}
\end{aligned}\end{equation*}
% RESIDUAL 85
\begin{equation*}\begin{aligned}
Q_{85}={}& x_{110}x_{203}-x_{111}x_{124}-x_{112}
\end{aligned}\end{equation*}
% RESIDUAL 86
\begin{equation*}\begin{aligned}
Q_{86}={}& x_{113}x_{204}-x_{115}x_{124}-x_{116}
\end{aligned}\end{equation*}
% RESIDUAL 87
\begin{equation*}\begin{aligned}
Q_{87}={}& x_{117}x_{204}-x_{118}x_{124}-x_{119}
\end{aligned}\end{equation*}
% RESIDUAL 88
\begin{equation*}\begin{aligned}
Q_{88}={}& -x_{18}x_{121}+x_{124}
\end{aligned}\end{equation*}
% RESIDUAL 89
\begin{equation*}\begin{aligned}
Q_{89}={}& x_{120}-x_{122}x_{124}-x_{123}
\end{aligned}\end{equation*}
% RESIDUAL 90
\begin{equation*}\begin{aligned}
Q_{90}={}& -x_{124}x_{126}+x_{125}-x_{127}
\end{aligned}\end{equation*}
% RESIDUAL 91
\begin{equation*}\begin{aligned}
Q_{91}={}& -x_{124}x_{129}+x_{128}x_{203}-x_{130}
\end{aligned}\end{equation*}
% RESIDUAL 92
\begin{equation*}\begin{aligned}
Q_{92}={}& -x_{124}x_{132}+x_{131}x_{203}-x_{133}
\end{aligned}\end{equation*}
% RESIDUAL 93
\begin{equation*}\begin{aligned}
Q_{93}={}& x_{134}-2x_{134}x_{136}+x_{137}
\end{aligned}\end{equation*}
% RESIDUAL 94
\begin{equation*}\begin{aligned}
Q_{94}={}& -x_{134}x_{137}-x_{137}+x_{138}
\end{aligned}\end{equation*}
% RESIDUAL 95
\begin{equation*}\begin{aligned}
Q_{95}={}& x_{135}+x_{135}^{2}+x_{137}^{2}-x_{138}^{2}-x_{189}x_{196}+x_{193}
\end{aligned}\end{equation*}
% RESIDUAL 96
\begin{equation*}\begin{aligned}
Q_{96}={}& -x_{135}x_{140}-x_{140}+x_{141}
\end{aligned}\end{equation*}
% RESIDUAL 97
\begin{equation*}\begin{aligned}
Q_{97}={}& -1+x_{139}^{2}+2x_{139}x_{141}+x_{140}^{2}
\end{aligned}\end{equation*}
% RESIDUAL 98
\begin{equation*}\begin{aligned}
Q_{98}={}& x_{3}^{2}-2x_{3}x_{135}-x_{135}+x_{142}+x_{189}x_{196}-x_{193}
\end{aligned}\end{equation*}
% RESIDUAL 99
\begin{equation*}\begin{aligned}
Q_{99}={}& -x_{3}x_{140}+x_{5}-x_{139}-x_{140}+x_{142}x_{143}
\end{aligned}\end{equation*}
% RESIDUAL 100
\begin{equation*}\begin{aligned}
Q_{100}={}& 1-2x_{4}+x_{4}^{2}-2x_{4}x_{135}+x_{144}
\end{aligned}\end{equation*}
% RESIDUAL 101
\begin{equation*}\begin{aligned}
Q_{101}={}& -x_{4}x_{176}+x_{18}+x_{144}x_{145}-x_{175}
\end{aligned}\end{equation*}
% RESIDUAL 102
\begin{equation*}\begin{aligned}
Q_{102}={}& -1+2x_{135}x_{187}-x_{146}+2x_{187}-x_{187}^{2}
\end{aligned}\end{equation*}
% RESIDUAL 103
\begin{equation*}\begin{aligned}
Q_{103}={}& -1+x_{95}+x_{146}x_{147}-x_{175}-x_{176}x_{187}
\end{aligned}\end{equation*}
% RESIDUAL 104
\begin{equation*}\begin{aligned}
Q_{104}={}& -x_{135}x_{149}-x_{149}+x_{150}
\end{aligned}\end{equation*}
% RESIDUAL 105
\begin{equation*}\begin{aligned}
Q_{105}={}& -1+x_{148}^{2}+2x_{148}x_{150}+x_{149}^{2}
\end{aligned}\end{equation*}
% RESIDUAL 106
\begin{equation*}\begin{aligned}
Q_{106}={}& 1+x_{135}-x_{135}x_{152}-x_{152}+x_{153}
\end{aligned}\end{equation*}
% RESIDUAL 107
\begin{equation*}\begin{aligned}
Q_{107}={}& x_{151}^{2}+2x_{151}x_{153}-2x_{152}+x_{152}^{2}
\end{aligned}\end{equation*}
% RESIDUAL 108
\begin{equation*}\begin{aligned}
Q_{108}={}& -x_{135}x_{155}-x_{155}+x_{156}
\end{aligned}\end{equation*}
% RESIDUAL 109
\begin{equation*}\begin{aligned}
Q_{109}={}& -1+x_{154}^{2}+2x_{154}x_{156}+x_{155}^{2}
\end{aligned}\end{equation*}
% RESIDUAL 110
\begin{equation*}\begin{aligned}
Q_{110}={}& -x_{135}x_{158}-x_{158}+x_{159}
\end{aligned}\end{equation*}
% RESIDUAL 111
\begin{equation*}\begin{aligned}
Q_{111}={}& -1+x_{157}^{2}+2x_{157}x_{159}+x_{158}^{2}
\end{aligned}\end{equation*}
% RESIDUAL 112
\begin{equation*}\begin{aligned}
Q_{112}={}& -x_{135}x_{161}-x_{161}+x_{162}
\end{aligned}\end{equation*}
% RESIDUAL 113
\begin{equation*}\begin{aligned}
Q_{113}={}& -1+x_{160}^{2}+2x_{160}x_{162}+x_{161}^{2}
\end{aligned}\end{equation*}
% RESIDUAL 114
\begin{equation*}\begin{aligned}
Q_{114}={}& -x_{135}x_{164}-x_{164}+x_{165}
\end{aligned}\end{equation*}
% RESIDUAL 115
\begin{equation*}\begin{aligned}
Q_{115}={}& -1+x_{163}^{2}+2x_{163}x_{165}+x_{164}^{2}
\end{aligned}\end{equation*}
% RESIDUAL 116
\begin{equation*}\begin{aligned}
Q_{116}={}& -x_{135}x_{167}-x_{167}+x_{168}
\end{aligned}\end{equation*}
% RESIDUAL 117
\begin{equation*}\begin{aligned}
Q_{117}={}& -1+x_{166}^{2}+2x_{166}x_{168}+x_{167}^{2}
\end{aligned}\end{equation*}
% RESIDUAL 118
\begin{equation*}\begin{aligned}
Q_{118}={}& -x_{135}x_{170}-x_{170}+x_{171}
\end{aligned}\end{equation*}
% RESIDUAL 119
\begin{equation*}\begin{aligned}
Q_{119}={}& -1+x_{169}^{2}+2x_{169}x_{171}+x_{170}^{2}
\end{aligned}\end{equation*}
% RESIDUAL 120
\begin{equation*}\begin{aligned}
Q_{120}={}& -x_{135}x_{173}-x_{173}+x_{174}
\end{aligned}\end{equation*}
% RESIDUAL 121
\begin{equation*}\begin{aligned}
Q_{121}={}& -1+x_{172}^{2}+2x_{172}x_{174}+x_{173}^{2}
\end{aligned}\end{equation*}
% RESIDUAL 122
\begin{equation*}\begin{aligned}
Q_{122}={}& 1-x_{135}x_{176}-x_{176}+x_{177}
\end{aligned}\end{equation*}
% RESIDUAL 123
\begin{equation*}\begin{aligned}
Q_{123}={}& -1+2x_{175}+x_{175}^{2}+2x_{175}x_{177}+x_{176}^{2}
\end{aligned}\end{equation*}
% RESIDUAL 124
\begin{equation*}\begin{aligned}
Q_{124}={}& -x_{135}x_{179}-x_{179}+x_{180}
\end{aligned}\end{equation*}
% RESIDUAL 125
\begin{equation*}\begin{aligned}
Q_{125}={}& -1+x_{178}^{2}+2x_{178}x_{180}+x_{179}^{2}
\end{aligned}\end{equation*}
% RESIDUAL 126
\begin{equation*}\begin{aligned}
Q_{126}={}& -x_{135}x_{182}-x_{182}+x_{183}
\end{aligned}\end{equation*}
% RESIDUAL 127
\begin{equation*}\begin{aligned}
Q_{127}={}& -1+x_{181}^{2}+2x_{181}x_{183}+x_{182}^{2}
\end{aligned}\end{equation*}
% RESIDUAL 128
\begin{equation*}\begin{aligned}
Q_{128}={}& -x_{135}x_{185}-x_{185}+x_{186}
\end{aligned}\end{equation*}
% RESIDUAL 129
\begin{equation*}\begin{aligned}
Q_{129}={}& -1+x_{184}^{2}+2x_{184}x_{186}+x_{185}^{2}
\end{aligned}\end{equation*}
% RESIDUAL 130
\begin{equation*}\begin{aligned}
Q_{130}={}& -2x_{0}-3x_{0}^{2}-2x_{0}x_{3}-x_{1}^{2}-x_{2}^{2}-x_{3}-2x_{3}^{2}-x_{4}^{2}-2x_{5}^{2}-2x_{11}\\
& +x_{11}^{2}+2x_{11}x_{20}-x_{12}^{2}-x_{13}^{2}-x_{14}^{2}-x_{15}^{2}-x_{16}^{2}-x_{17}^{2}-x_{18}^{2}\\
& -x_{24}-x_{25}+x_{95}-2x_{95}^{2}-x_{96}^{2}-x_{97}^{2}-x_{101}^{2}-x_{104}^{2}-x_{107}^{2}-x_{110}^{2}\\
& -x_{113}^{2}-x_{114}^{2}-x_{117}^{2}-x_{120}^{2}-x_{121}^{2}-x_{125}^{2}-x_{128}^{2}-x_{131}^{2}+x_{134}\\
& +x_{140}^{2}-x_{146}x_{147}+x_{149}^{2}-2x_{152}+x_{152}^{2}+x_{155}^{2}+x_{158}^{2}+x_{161}^{2}\\
& +x_{164}^{2}+x_{167}^{2}+x_{170}^{2}+x_{173}^{2}+x_{175}+x_{176}^{2}+x_{176}x_{187}+x_{179}^{2}\\
& +x_{182}^{2}-x_{185}^{2}-x_{187}^{2}+x_{205}^{2}+x_{206}^{2}+x_{207}+x_{207}^{2}
\end{aligned}\end{equation*}
% RESIDUAL 131
\begin{equation*}\begin{aligned}
Q_{131}={}& x_{4}-2x_{4}x_{11}-2x_{4}x_{12}-2x_{4}x_{13}-2x_{4}x_{14}-2x_{4}x_{15}-2x_{4}x_{16}\\
& -2x_{4}x_{17}-2x_{4}x_{20}+x_{187}
\end{aligned}\end{equation*}
% RESIDUAL 132
\begin{equation*}\begin{aligned}
Q_{132}={}& -2+x_{0}-2x_{1}+x_{2}+x_{3}-x_{4}x_{19}+x_{5}x_{6}+x_{7}x_{8}-2x_{10}-2x_{10}x_{11}\\
& -x_{10}x_{20}+x_{11}+2x_{11}^{2}-3x_{11}x_{20}+x_{17}+x_{20}+2x_{20}^{2}-3x_{21}-x_{24}\\
& -x_{72}-x_{135}x_{176}-x_{176}+x_{177}-x_{188}^{2}-x_{189}^{2}-x_{190}^{2}-x_{191}^{2}-x_{192}^{2}\\
& +x_{199}^{2}+x_{200}^{2}
\end{aligned}\end{equation*}
% RESIDUAL 133
\begin{equation*}\begin{aligned}
Q_{133}={}& -x_{135}x_{190}-x_{190}+x_{199}
\end{aligned}\end{equation*}
% RESIDUAL 134
\begin{equation*}\begin{aligned}
Q_{134}={}& -1+x_{189}^{2}+x_{190}^{2}-x_{199}^{2}
\end{aligned}\end{equation*}
% RESIDUAL 135
\begin{equation*}\begin{aligned}
Q_{135}={}& -x_{192}-x_{192}x_{193}+x_{200}
\end{aligned}\end{equation*}
% RESIDUAL 136
\begin{equation*}\begin{aligned}
Q_{136}={}& -1+x_{191}^{2}+x_{192}^{2}-x_{200}^{2}
\end{aligned}\end{equation*}
% RESIDUAL 137
\begin{equation*}\begin{aligned}
Q_{137}={}& -x_{185}^{2}+x_{201}
\end{aligned}\end{equation*}
% RESIDUAL 138
\begin{equation*}\begin{aligned}
Q_{138}={}& x_{190}-2x_{195}x_{201}-x_{201}
\end{aligned}\end{equation*}
% RESIDUAL 139
\begin{equation*}\begin{aligned}
Q_{139}={}& x_{135}x_{190}-3x_{190}-8x_{190}x_{194}+x_{193}-x_{199}
\end{aligned}\end{equation*}
% RESIDUAL 140
\begin{equation*}\begin{aligned}
Q_{140}={}& -x_{135}-x_{189}x_{196}+x_{193}
\end{aligned}\end{equation*}
% RESIDUAL 141
\begin{equation*}\begin{aligned}
Q_{141}={}& -x_{184}-x_{186}-x_{189}x_{197}+x_{191}
\end{aligned}\end{equation*}
% RESIDUAL 142
\begin{equation*}\begin{aligned}
Q_{142}={}& -2x_{10}-x_{20}-2x_{20}^{2}-2x_{21}-2x_{23}-2x_{27}-2x_{29}-2x_{31}-2x_{33}-4x_{35}\\
& -2x_{36}-2x_{39}-2x_{42}-4x_{61}-2x_{62}-4x_{185}x_{198}+x_{192}
\end{aligned}\end{equation*}
% RESIDUAL 143
\begin{equation*}\begin{aligned}
Q_{143}={}& 2-x_{0}-x_{0}^{2}-2x_{0}x_{140}-2x_{0}x_{149}+x_{1}-2x_{1}x_{176}-x_{2}^{2}-x_{3}\\
& -x_{3}^{2}+x_{3}x_{140}-2x_{3}x_{149}-x_{4}^{2}-x_{5}-x_{5}^{2}-2x_{5}x_{167}-x_{7}x_{8}\\
& +2x_{10}-2x_{10}x_{152}+2x_{10}x_{155}-8x_{10}x_{185}-2x_{11}x_{158}-2x_{12}x_{161}\\
& -2x_{13}x_{164}-2x_{14}x_{170}-2x_{15}x_{173}-2x_{16}x_{179}-2x_{17}x_{182}-x_{18}^{2}\\
& +4x_{20}-4x_{20}x_{152}+2x_{20}x_{155}-2x_{20}x_{185}-8x_{21}x_{185}-4x_{23}x_{185}\\
& -4x_{27}x_{185}-4x_{29}x_{185}-4x_{31}x_{185}-4x_{33}x_{185}-8x_{35}x_{185}-4x_{36}x_{185}\\
& -4x_{39}x_{185}-4x_{42}x_{185}-8x_{61}x_{185}-4x_{62}x_{185}+2x_{95}-2x_{95}^{2}-x_{96}^{2}\\
& -x_{97}^{2}-x_{101}^{2}-x_{104}^{2}-x_{107}^{2}-x_{110}^{2}-x_{113}^{2}-x_{114}^{2}-x_{117}^{2}-x_{120}^{2}\\
& -x_{121}^{2}-x_{125}^{2}-x_{128}^{2}-x_{131}^{2}+x_{134}-x_{135}x_{140}-x_{135}x_{149}\\
& -x_{135}x_{155}-x_{135}x_{158}-x_{135}x_{161}-x_{135}x_{164}-x_{135}x_{167}-x_{135}x_{170}\\
& -x_{135}x_{173}-x_{135}x_{179}-x_{135}x_{182}-3x_{135}x_{185}-x_{135}x_{202}-x_{139}-x_{141}\\
& +x_{141}^{2}-x_{142}x_{143}-2x_{148}-x_{149}-x_{150}+x_{150}^{2}-2x_{151}+2x_{152}-2x_{153}\\
& +x_{153}^{2}-2x_{154}-x_{155}-x_{156}+x_{156}^{2}-2x_{157}-x_{158}-x_{159}+x_{159}^{2}\\
& -2x_{160}-x_{161}-x_{162}+x_{162}^{2}-2x_{163}-x_{164}-x_{165}+x_{165}^{2}-2x_{166}-x_{167}\\
& -x_{168}+x_{168}^{2}-2x_{169}-x_{170}-x_{171}+x_{171}^{2}-2x_{172}-x_{173}-x_{174}+x_{174}^{2}\\
& +x_{175}^{2}+2x_{175}x_{177}+x_{176}^{2}+x_{177}^{2}-2x_{178}-x_{179}-x_{180}+x_{180}^{2}\\
& -2x_{181}-x_{182}-x_{183}+x_{183}^{2}-2x_{184}+5x_{185}+x_{186}+x_{186}^{2}-x_{187}^{2}\\
& -x_{188}^{2}
\end{aligned}\end{equation*}
% RESIDUAL 144
\begin{equation*}\begin{aligned}
Q_{144}={}& -x_{97}x_{203}+x_{121}
\end{aligned}\end{equation*}
% RESIDUAL 145
\begin{equation*}\begin{aligned}
Q_{145}={}& -x_{114}x_{204}+x_{121}
\end{aligned}\end{equation*}
% RESIDUAL 146
\begin{equation*}\begin{aligned}
Q_{146}={}& 1-5x_{10}+2x_{10}^{2}-2x_{10}x_{11}-x_{10}x_{20}-16x_{10}x_{23}-16x_{10}x_{27}\\
& -16x_{10}x_{29}-16x_{10}x_{31}-16x_{10}x_{33}-24x_{10}x_{35}-16x_{10}x_{36}\\
& -16x_{10}x_{39}-16x_{10}x_{42}-32x_{10}x_{61}-16x_{10}x_{62}-3x_{11}+2x_{11}^{2}\\
& +x_{11}x_{20}-3x_{20}-5x_{20}^{2}-8x_{20}x_{21}-4x_{20}x_{23}-4x_{20}x_{27}-4x_{20}x_{29}\\
& -4x_{20}x_{31}-4x_{20}x_{33}-4x_{20}x_{36}-4x_{20}x_{39}-4x_{20}x_{42}-8x_{20}x_{61}\\
& -4x_{20}x_{62}-16x_{21}x_{23}-16x_{21}x_{27}-16x_{21}x_{29}-16x_{21}x_{31}\\
& -16x_{21}x_{33}-32x_{21}x_{35}-16x_{21}x_{36}-16x_{21}x_{39}-16x_{21}x_{42}\\
& -32x_{21}x_{61}-16x_{21}x_{62}+4x_{23}-4x_{23}^{2}-8x_{23}x_{27}-8x_{23}x_{29}\\
& -8x_{23}x_{31}-8x_{23}x_{33}-16x_{23}x_{35}-8x_{23}x_{36}-8x_{23}x_{39}-8x_{23}x_{42}\\
& -16x_{23}x_{61}-8x_{23}x_{62}-2x_{24}-x_{25}+4x_{27}-4x_{27}^{2}-8x_{27}x_{29}\\
& -8x_{27}x_{31}-8x_{27}x_{33}-16x_{27}x_{35}-8x_{27}x_{36}-8x_{27}x_{39}-8x_{27}x_{42}\\
& -16x_{27}x_{61}-8x_{27}x_{62}+4x_{29}-4x_{29}^{2}-8x_{29}x_{31}-8x_{29}x_{33}\\
& -16x_{29}x_{35}-8x_{29}x_{36}-8x_{29}x_{39}-8x_{29}x_{42}-16x_{29}x_{61}-8x_{29}x_{62}\\
& +4x_{31}-4x_{31}^{2}-8x_{31}x_{33}-16x_{31}x_{35}-8x_{31}x_{36}-8x_{31}x_{39}\\
& -8x_{31}x_{42}-16x_{31}x_{61}-8x_{31}x_{62}+4x_{33}-4x_{33}^{2}-16x_{33}x_{35}\\
& -8x_{33}x_{36}-8x_{33}x_{39}-8x_{33}x_{42}-16x_{33}x_{61}-8x_{33}x_{62}+8x_{35}\\
& -24x_{35}^{2}-16x_{35}x_{36}-16x_{35}x_{39}-16x_{35}x_{42}-32x_{35}x_{61}-16x_{35}x_{62}\\
& +6x_{36}-4x_{36}^{2}-8x_{36}x_{39}-8x_{36}x_{42}-16x_{36}x_{61}-8x_{36}x_{62}-x_{37}\\
& +4x_{39}-4x_{39}^{2}-8x_{39}x_{42}-16x_{39}x_{61}-8x_{39}x_{62}+4x_{42}-4x_{42}^{2}\\
& -16x_{42}x_{61}-8x_{42}x_{62}+12x_{61}-12x_{61}^{2}-16x_{61}x_{62}+3x_{62}-4x_{62}^{2}\\
& +x_{72}-x_{140}^{2}-x_{149}^{2}+2x_{152}-x_{152}^{2}-x_{155}^{2}-x_{158}^{2}-x_{161}^{2}-x_{164}^{2}\\
& -x_{167}^{2}-x_{170}^{2}-x_{173}^{2}-x_{176}^{2}-x_{179}^{2}-x_{182}^{2}+x_{185}^{2}-24x_{185}x_{198}\\
& +4x_{192}-2x_{192}x_{193}+2x_{200}-x_{205}^{2}-x_{206}^{2}-x_{207}-x_{207}^{2}
\end{aligned}\end{equation*}
% RESIDUAL 147
\begin{equation*}\begin{aligned}
Q_{147}={}& -1-4x_{99}^{2}-4x_{103}^{2}-4x_{106}^{2}-4x_{109}^{2}-4x_{112}^{2}-4x_{116}^{2}-4x_{119}^{2}+x_{121}^{2}\\
& -4x_{123}^{2}-4x_{127}^{2}-4x_{130}^{2}-4x_{133}^{2}-x_{208}^{2}-x_{209}^{2}-x_{210}-x_{210}^{2}\\
& -4x_{212}^{2}
\end{aligned}\end{equation*}
% RESIDUAL 148
\begin{equation*}\begin{aligned}
Q_{148}={}& -x_{4}x_{211}+x_{121}
\end{aligned}\end{equation*}
% RESIDUAL 149
\begin{equation*}\begin{aligned}
Q_{149}={}& x_{2}x_{121}-x_{5}x_{211}-x_{212}
\end{aligned}\end{equation*}
% RESIDUAL 150
\begin{equation*}\begin{aligned}
Q_{150}={}& -x_{2}x_{13}+x_{213}
\end{aligned}\end{equation*}
% RESIDUAL 151
\begin{equation*}\begin{aligned}
Q_{151}={}& x_{95}x_{215}-2x_{96}x_{149}+2x_{101}-2x_{101}x_{152}-2x_{104}x_{155}-2x_{107}x_{158}\\
& -2x_{110}x_{161}-2x_{113}x_{164}-2x_{117}x_{170}-2x_{120}x_{173}-2x_{125}x_{179}\\
& -2x_{128}x_{182}-2x_{131}x_{185}-2x_{135}x_{215}+2x_{148}x_{149}+2x_{149}x_{150}-2x_{151}\\
& +2x_{151}x_{152}+2x_{152}x_{153}-2x_{153}+2x_{154}x_{155}+2x_{155}x_{156}+2x_{157}x_{158}\\
& +2x_{158}x_{159}+2x_{160}x_{161}+2x_{161}x_{162}+2x_{163}x_{164}+2x_{164}x_{165}\\
& +2x_{169}x_{170}+2x_{170}x_{171}+2x_{172}x_{173}+2x_{173}x_{174}+2x_{178}x_{179}\\
& +2x_{179}x_{180}+2x_{181}x_{182}+2x_{182}x_{183}+2x_{184}x_{185}+2x_{185}x_{186}+x_{214}\\
& -2x_{215}
\end{aligned}\end{equation*}
% RESIDUAL 152
\begin{equation*}\begin{aligned}
Q_{152}={}& -1-x_{95}x_{214}+x_{96}^{2}-2x_{96}x_{148}+x_{101}^{2}-2x_{101}x_{151}+x_{104}^{2}\\
& -2x_{104}x_{154}+x_{107}^{2}-2x_{107}x_{157}+x_{110}^{2}-2x_{110}x_{160}+x_{113}^{2}\\
& -2x_{113}x_{163}+x_{117}^{2}-2x_{117}x_{169}+x_{120}^{2}-2x_{120}x_{172}+x_{125}^{2}\\
& -2x_{125}x_{178}+x_{128}^{2}-2x_{128}x_{181}+x_{131}^{2}-2x_{131}x_{184}+x_{148}^{2}-x_{149}^{2}\\
& +x_{151}^{2}+2x_{152}-x_{152}^{2}+x_{154}^{2}-x_{155}^{2}+x_{157}^{2}-x_{158}^{2}+x_{160}^{2}-x_{161}^{2}\\
& +x_{163}^{2}-x_{164}^{2}+x_{169}^{2}-x_{170}^{2}+x_{172}^{2}-x_{173}^{2}+x_{178}^{2}-x_{179}^{2}+x_{181}^{2}\\
& -x_{182}^{2}+x_{184}^{2}-x_{185}^{2}+x_{215}
\end{aligned}\end{equation*}
% RESIDUAL 153
\begin{equation*}\begin{aligned}
Q_{153}={}& x_{95}x_{217}-2x_{97}x_{140}-2x_{114}x_{167}-2x_{121}x_{176}-2x_{135}x_{217}\\
& +2x_{139}x_{140}+2x_{140}x_{141}+2x_{166}x_{167}+2x_{167}x_{168}+2x_{175}x_{176}+2x_{176}\\
& +2x_{176}x_{177}+x_{216}-3x_{217}
\end{aligned}\end{equation*}
% RESIDUAL 154
\begin{equation*}\begin{aligned}
Q_{154}={}& -x_{95}x_{216}+x_{97}^{2}-2x_{97}x_{139}+x_{114}^{2}-2x_{114}x_{166}+x_{121}^{2}-2x_{121}x_{175}\\
& +x_{139}^{2}-x_{140}^{2}+x_{166}^{2}-x_{167}^{2}+x_{175}^{2}-x_{176}^{2}+x_{216}+x_{217}
\end{aligned}\end{equation*}

\endgroup

\section{Exact data conventions and identities}

A sparse term \code{[c,[i,j,...]]} denotes
$c\,x_i x_j\cdots$; an empty index list denotes the constant $c$.
Repeated indices are repeated factors. A residual file declares its
coordinate order, unrestricted-integer domain, and sum-of-squares
combiner. Expanded files list the coefficients of the single collected
polynomial itself. SHA-256 identifiers authenticate bytes or the named
canonical JSON encoding; they do not replace the polynomial definitions.

\begin{longtable}{@{}p{0.26\textwidth}p{0.69\textwidth}@{}}
\toprule Artifact & SHA-256 of the supplied file\\\midrule
quartic218/residuals & {\scriptsize\ttfamily 878d2f427a117f35b4c211a223003978\newline d52509acc46a5c97b829380b6e608ad8}\\
quartic218/expanded polynomial & {\scriptsize\ttfamily f81c0390e593d503d9274bd20c718385\newline f8a20b6f7aa3f2fa5b74a43cf0157c34}\\
decic215/residuals & {\scriptsize\ttfamily a335e0d1d1d82aa2cd30c2f436ed61a7\newline 52fff001db134b339c23fae01d4a70b0}\\
decic215/expanded polynomial & {\scriptsize\ttfamily 3f2c8eadf41cc85b54b5a55ad1cfb4fd\newline 0f5100f1373723c706e836a788737d6e}\\
\bottomrule\end{longtable}


The frozen transformation plan starts at the retained original
225-variable residual file, replays five radical projections, the affine
hyperplane chart, one affine substitution, and the protected row basis,
then emits the reciprocal companion. The final standalone equations can
be evaluated without executing a source generator or opening an archive.

\clearpage
\begin{thebibliography}{9}\small
\bibitem{parent158}
V. Reshetnikov's ProveIt research collection, AI-assisted contribution,
\emph{Pooling exact power certificates for Diophantine consistency},
package 158, 25 October 2026.
\url{https://github.com/VladimirReshetnikov/ProveIt/tree/953fb4d82d4229f4efc5f9db4cc2bf38ae329bb4/Computability/DiophantineConsistency/con-zfc-diophantine-forms/158-pooled-powers}.

\bibitem{projection159}
ProveIt research collection, AI-assisted contribution,
\emph{A direct quartic for Con(ZFC) in 234 integer variables}, package 159,
25 October 2026. Exact projection lemmas and source utilities.
\url{https://github.com/VladimirReshetnikov/ProveIt/tree/953fb4d82d4229f4efc5f9db4cc2bf38ae329bb4/Computability/DiophantineConsistency/con-zfc-diophantine-forms/159-quadratic-projection-234}.

\bibitem{catalog}
ProveIt, \emph{Con(ZFC) in Diophantine form}, collection README at the
comparison snapshot. The summary frontier predates some latest leaf packages.
\url{https://github.com/VladimirReshetnikov/ProveIt/blob/953fb4d82d4229f4efc5f9db4cc2bf38ae329bb4/Computability/DiophantineConsistency/con-zfc-diophantine-forms/README.md}.

\bibitem{current}
ProveIt, \emph{A Diophantine equation for syntactic consistency of ZFC},
active research README and complete displayed equation, comparison snapshot.
\url{https://github.com/VladimirReshetnikov/ProveIt/tree/953fb4d82d4229f4efc5f9db4cc2bf38ae329bb4/Computability/DiophantineConsistency}.

\bibitem{reciprocal660}
ProveIt research collection, \emph{Square-shell syntax and
reciprocal-factorial CRT}, preserved under package 112.
\url{https://github.com/VladimirReshetnikov/ProveIt/tree/953fb4d82d4229f4efc5f9db4cc2bf38ae329bb4/Computability/DiophantineConsistency/con-zfc-diophantine-forms/112-powered-saturation-round14}.

\bibitem{mr75}
Yu. Matiyasevich and J. Robinson,
\emph{Reduction of an arbitrary Diophantine equation to one in 13 unknowns},
Acta Arithmetica 27 (1975), 521--553.
\url{https://doi.org/10.4064/aa-27-1-521-553}.

\bibitem{bayer2025}
J. Bayer, M. David, M. Ha\ss ler, Yu. Matiyasevich, and D. Schleicher,
\emph{Diophantine equations over $\Z$: universal bounds and parallel
formalization}, arXiv:2506.20909v1 (2025).
\url{https://arxiv.org/html/2506.20909v1}.
\end{thebibliography}
\end{document}
